% !TEX program = xelatex
% !TeX encoding = UTF-8

\documentclass[11pt,a4paper]{article}

% ============================================================
% 한국어 설정 (XeLaTeX + xeCJK)
% ============================================================
\usepackage{fontspec}
\setmainfont{Liberation Serif}          % 라틴 문자용
\usepackage{xeCJK}
\setCJKmainfont{Noto Sans KR}           % 한국어 메인 (시스템 폰트)
\setCJKmonofont{Noto Sans KR}           % 고정폭
\XeTeXlinebreaklocale "ko"              % 한국어 줄바꿈
\XeTeXlinebreakskip = 0pt plus 1pt

% ============================================================
% 수학 및 레이아웃
% ============================================================
\usepackage{amsmath,amssymb,amsthm,bm}
\usepackage{geometry}
\usepackage{mathtools}
\usepackage{enumitem}
\usepackage{siunitx}
\usepackage{booktabs}
\usepackage{array}
\usepackage{slashed}
\usepackage{graphicx}
\usepackage{caption}
\usepackage{subcaption}
\usepackage{braket}
\usepackage{tensor}
\usepackage{makecell}
\usepackage[most]{tcolorbox}
\usepackage{pgfplots}
\usepackage{float}
\usepackage{tikz}
\usepackage{multicol}
\usepackage{lipsum}
\usepackage{microtype}
\usepackage{hyperref}

% ============================================================
% 누락 패키지 및 색상 정의 (컴파일 오류 방지)
% ============================================================
\usepackage{longtable}
\usepackage{xcolor}
\definecolor{darkblue}{RGB}{0,0,139}

\pgfplotsset{compat=1.18}
\usetikzlibrary{arrows.meta,positioning,shapes.geometric,fit,calc,
	decorations.pathreplacing,patterns,quotes,angles,
	shadows,decorations.markings}

\geometry{margin=1in}
\setlength{\emergencystretch}{3em}
\sloppy

% ============================================================
% 리스트 설정 (Python 코드용)
% ============================================================
\usepackage{listings}
\lstset{
	language=Python,
	basicstyle=\ttfamily\small,
	breaklines=true,
	showstringspaces=false,
	frame=single,
	columns=flexible,
	extendedchars=true,
}

% ============================================================
% 정리 환경 (한국어)
% ============================================================
\newtheorem{theorem}{정리}[section]
\newtheorem{lemma}{보조정리}[section]
\newtheorem{axiom}{공리}[section]
\newtheorem{remark}{주석}[section]
\newtheorem{proposition}{명제}[section]
\newtheorem{definition}{정의}[section]
\newtheorem{assumption}{가정}[section]
\newtheorem{corollary}{따름정리}[section]
\newtheorem{example}{예시}[section]

% ============================================================
% YUCT 명령어 (변경 없음)
% ============================================================
\NewDocumentCommand{\Keff}{o}{%
	K_{\mathrm{eff}\IfValueT{#1}{,#1}}%
}
\newcommand{\kappac}{\kappa_c}
\newcommand{\eps}{\varepsilon}
\newcommand{\betaf}{\beta}
\newcommand{\Sodd}{S_{\mathrm{odd}}}
\newcommand{\Seven}{S_{\mathrm{even}}}
\newcommand{\calD}{\mathcal{D}}
\newcommand{\calI}{\mathcal{I}}
\newcommand{\calR}{\mathcal{R}}
\newcommand{\Mpl}{M_{\mathrm{Pl}}}
\newcommand{\Lpl}{\ell_{\mathrm{Pl}}}
\newcommand{\tPl}{t_{\mathrm{Pl}}}
\newcommand{\Rcrit}{R_{\mathrm{crit}}}
\newcommand{\Cpers}{\mathbf{C}_{\text{pers}}}
\newcommand{\Yak}{\mathrm{Yak}}

% ============================================================
% PDF 메타데이터
% ============================================================
\hypersetup{
	pdftitle={철학의 정량적 기술: 형이상학에서 조정 물리학으로},
	pdfauthor={Alexey V. Yakushev},
	pdfsubject={철학의 정량적 분석, 조정 효율, D+I•R 삼원소, 일반화 벨 부등식},
	pdfkeywords={YUCT, YPSDC, 조정 효율, 철학, 정량적 철학, 의미 분석, 벨 부등식},
	breaklinks=true,
	pdfencoding=auto,
	bookmarksnumbered=true,
	bookmarksopen=true,
	bookmarksopenlevel=1
}

% ============================================================
% 본문 시작
% ============================================================
\begin{document}
	
	% ============================================================
	% 표지
	% ============================================================
	\begin{titlepage}
		\begin{center}
			\vspace*{1cm}
			{\Huge\bfseries\color{darkblue} 철학의 정량적 기술:\par}
			\vspace{0.3cm}
			{\Huge\bfseries\color{darkblue} 형이상학에서 조정 물리학으로\par}
			\vspace{1cm}
			{\Large\bfseries 인류 역사상 최초의 엄밀한 철학 체계의 정량적 기술\par}
			\vspace{1.5cm}
			{\large Alexey V. Yakushev\par}
			{\large Yakushev Research\par}
			{\large \url{https://yuct.org/}\par}
			\vspace{1cm}
			{\large \textbf{DOI:} \href{https://doi.org/10.5281/zenodo.18444599}{10.5281/zenodo.18444599}\par}
			\vspace{1cm}
			{\large 버전 1.3.1 \quad 2026년 6월\par}
			\vfill
			\begin{quote}
				\itshape
				「철학은 더 이상 추론의 예술이 아니라, 물리학과 동일한 법칙을 따르는 엄격한 학문이 된다.」
			\end{quote}
		\end{center}
	\end{titlepage}
	
	% ============================================================
	% 목차
	% ============================================================
	\tableofcontents
	\newpage
	
	% ============================================================
	% 서론
	% ============================================================
	\section{서론: 철학의 역설}
	
	인류 역사를 통해 철학은 「학문의 여왕」—존재, 인식, 진리, 의미에 대해 질문하는 학문—이었다. 그러나 물리학, 화학, 생물학과 달리 철학은 결코 \emph{정량적} 언어를 가지지 못했다. 그것은 해석의 예술이었고, 끝없는 논쟁의 장이었으며, 각 학파는 스스로를 「심오하다」고 주장할 수 있었지만 누구도 그 깊이를 측정할 수 없었다.
	
	이 역설—「철학은 모든 것을 말하지만 아무것도 측정할 수 없다」—은 수세기 동안 지속되었다. 칸트는 인식의 구조를 기술했지만 범주에 대한 척도를 제공하지 않았다. 헤겔은 변증법을 구축했지만 그 방정식을 도출하지 않았다. 포퍼는 경계 설정 기준(반증 가능성)을 제안했지만 \emph{반증 가능성의 정도}를 측정하지 않았다. 섀넌은 정보의 \emph{양}을 측정했지만 그 \emph{깊이}나 철학 체계의 \emph{연계성}을 측정하지 않았다.
	
	\begin{quote}
		\textit{「철학은 모든 것을 알지만 그것을 측정하는 방법을 모르는 과학이다.」}
	\end{quote}
	
	2026년, 이 역설은 해결되었다. 통합 조정 이론(YUCT)은 철학 체계의 정량적 기술을 위한 \textbf{엄밀한 수학적 장치}를 최초로 제공한다.
	
	이 돌파구의 핵심에는 단순하지만 심오한 관찰이 있다:
	
	\begin{center}
		\fbox{\parbox{0.8\textwidth}{\centering\Large\bfseries
				어떤 철학 체계든 조정이다.\\
				조정은 측정 가능하다.\\
				따라서 철학은 측정 가능하다.
		}}
	\end{center}
	
	본고는 인류 역사상 최초의 체계적인 철학의 정량적 기술을 제시한다. 다음을 보여줄 것이다:
	
	\begin{enumerate}
		\item 어떤 철학 체계든 삼원소 \(D + I \cdot R\)로 표현될 수 있다.
		\item 철학 체계의 깊이는 조정 효율 \(\Keff\)로 측정된다.
		\item 내부 모순은 조정 오차 \(\eps\)로 측정된다.
		\item 존재론적 연계성은 일반화 벨 부등식 \(S(\Keff)\)로 측정된다.
		\item 철학사는 \(\Keff\)의 진화이다.
		\item 영혼은 철학 범주로서 엄밀한 정량적 정의를 얻는다.
	\end{enumerate}
	
	우리는 「또 다른 철학 이론」을 제안하는 것이 아니다. 우리가 제안하는 것은 \textbf{철학의 물리학}—생각이 현실 및 자기 자신과 어떻게 조정되는지에 관한 엄밀한 과학이다.
	
	\section{역사-방법론적 새로움의 논증}
	
	「인류 역사상 최초의 엄밀한 철학 체계의 정량적 기술」이라는 주장은 과장된 것처럼 들릴 수 있다. 그러나 과학사와 철학사의 건조한 관점에서 보면, 그것은 엄격한 역사적 사실의 진술이다.
	
	과거의 위대한 사상가들이 철학이나 의식을 디지털화하려는 시도는, YUCT가 최초로 제공하는 다음 세 가지 핵심 요소가 결여되었기 때문에 방법론적 막다른 골목에 빠졌다:
	
	\begin{enumerate}
		\item \textbf{정보 이론}(섀넌, 1948)—정보의 \emph{양}을 측정할 수 있었지만 그 \emph{의미}는 측정할 수 없었다.
		\item \textbf{복잡성 이론}과 \textbf{프랙탈}(만델브로, 1975)—자기 유사 구조를 기술할 수 있었지만 그들의 조정 효율은 기술할 수 없었다.
		\item \textbf{조정}이라는 존재론적 원초 개념—완전히 결여되어 있었다.
	\end{enumerate}
	
	\subsection{왜 이전 시도들은 실패했는가}
	
	\begin{table}[H]
		\centering
		\caption{사고의 정량적 기술 방법론 비교}
		\label{tab:method-comparison}
		\begin{tabular}{p{3cm}p{3.5cm}p{3.5cm}p{4cm}}
			\toprule
			\textbf{사상가/학파} & \textbf{제안} & \textbf{막다른 골목} & \textbf{결여된 것} \\
			\midrule
			피타고라스 & 「만물은 수」 & 신비적 기하학; 관념의 의미 측정 불가 & 정보 이론, 엔트로피 \\
			라이프니츠 & \emph{보편 문자} + \emph{계산하라!} & 무엇을 계산해야 할지 불명 & 엔트로피와 사전 데이터베이스 \\
			섀넌 & 정보 엔트로피 \( H = -\sum p \log p \) & 의미를 의도적으로 배제 & 공명 \(R\)과 조정 효율 \(\Keff\) \\
			\bottomrule
		\end{tabular}
	\end{table}
	
	물론, 이 사상가들은 과학과 철학의 발전에 중요한 기여를 했다. YUCT는 그들의 업적을 부정하지 않고 체계화하고 보완하며, 그들에게 부족했던 도구를 제공한다. 이전 시도들은 「오류」가 아니라 오늘날의 과제에 대해 단순히 \emph{불완전}했던 것이다.
	
	\subsection{YUCT가 원칙적으로 새롭게 추가하는 것}
	
	모든 선배들과 달리 YUCT는 다음을 제공한다:
	
	\begin{enumerate}
		\item \textbf{조정 효율의 척도:}
		\[
		\Keff = \frac{H(D)}{H(I)}
		\]
		여기서 \(H(D)\)는 사전(모든 개념)의 엔트로피, \(H(I)\)는 인덱스(기본 공리)의 엔트로피이다.
		
		\item \textbf{철학적 오차의 공식:}
		\[
		\eps_{\text{철}} = \kappac \cdot \alpha_{\text{철}} \cdot \Keff^{-2/3}
		\]
		이는 기본 입자의 질량과 동일한 플랑크 불변량에서 유도된다.
		
		\item \textbf{철학을 위한 일반화 벨 부등식:}
		\[
		S(\Keff) = 2 + (2\sqrt{2} - 2) \cdot \left(1 - 2\kappac \alpha \Keff^{-2/3}\right)
		\]
		이는 체계의 존재론적 연계성을 측정한다.
		
		\item \textbf{검증 가능성:} 이 모든 양은 텍스트의 의미 분석을 통해 계산 가능하며, 철학을 기기 검증 가능한 학문으로 만든다.
	\end{enumerate}
	
	\section{삼원소 \texorpdfstring{\(D + I \cdot R\)}{D + I * R}: 존재론적 원초 개념}
	
	\subsection{정의}
	
	YUCT에서는 모든 시스템—물리적, 생물학적, 정보적, 철학적—이 삼원소로 기술된다:
	
	\begin{equation}
		\boxed{\text{실재} = D + I \cdot R}
	\end{equation}
	
	여기서:
	
	\begin{definition}[사전 \(D\)]
		\textbf{사전}—시스템이 활성화할 수 있는 모든 가능한 상태, 개념, 규칙, 판단의 집합. 철학에서는 범주, 공리, 기본 원리의 체계이다.
	\end{definition}
	
	\begin{definition}[인덱스 \(I\)]
		\textbf{인덱스}—현재 상태, 사전의 특정 요소에 대한 지시. 철학에서는 기본 공리, 즉 전체 체계를 전개하는 근본적 명제이다.
	\end{definition}
	
	\begin{definition}[공명 \(R\)]
		\textbf{공명}—인덱스가 사전을 얼마나 완전하고 안정적으로 활성화하는지의 척도. 철학에서는 체계의 연계성, 즉 하나의 공리가 전체 존재론을 결정하는 정도이다.
	\end{definition}
	
	\subsection{조정으로서의 철학 체계}
	
	어떤 철학 체계든 조정의 특수한 예이다:
	
	\begin{itemize}
		\item \textbf{사전} \(D\)—철학자가 사용하는 모든 개념(범주, 판단, 용어).
		\item \textbf{인덱스} \(I\)—모든 것을 유도하는 기본 공리(예: 데카르트의 \textit{나는 생각한다, 고로 존재한다}).
		\item \textbf{공명} \(R\)—이 공리들이 체계 전체를 결정하는 정도(하나의 명제가 철학 전체를 「깊이」 관통하는 정도).
	\end{itemize}
	
	\begin{example}[데카르트 철학]
		\begin{itemize}
			\item \(D\)—데카르트 철학의 모든 개념: 실체, 사유, 연장, 신, 의심 등.
			\item \(I\)—\textit{나는 생각한다, 고로 존재한다}.
			\item \(R\)—이 단일 인덱스에서 데카르트 형이상학 전체가 전개된다: 이원론, 신 존재 증명, 인식의 본질.
		\end{itemize}
	\end{example}
	
	\section{철학의 조정 효율 \texorpdfstring{\(\Keff\)}{K\_eff}}
	
	\subsection{정의}
	
	철학 체계의 조정 효율은 사전 엔트로피와 인덱스 엔트로피의 비율로 정의된다:
	
	\begin{equation}
		\boxed{
			\Keff[\text{철}] = \frac{H(D)}{H(I)}
		}
		\label{eq:keff-phil}
	\end{equation}
	
	여기서 \(H(D)\)는 사전의 정보 용량(구별되는 개념 수), \(H(I)\)는 인덱스의 정보 용량(기본 공리 수)이다.
	
	\textbf{해석:} \(\Keff[\text{철}]\)은 \emph{최소 수의 공리에서 얼마나 많은 의미가 생성되는지}를 측정한다. \(\Keff\)가 높을수록 철학 체계는 더 「깊고」 「포괄적」이다.
	
	\subsection{사전과 인덱스의 측정}
	
	YUCT에서 사전과 인덱스는 다음을 통해 측정 가능하다:
	
	\begin{enumerate}
		\item \textbf{텍스트의 의미 분석:} 의미 네트워크의 크기(\(H(D)\))는 고유 개념과 그 연결의 수로 측정된다.
		\item \textbf{공리계 분석:} 인덱스의 길이(\(H(I)\))는 체계를 유도하는 근본적 명제의 수로 측정된다.
		\item \textbf{연계성:} 공명 \(R\)은 의미 네트워크 내의 평균 경로 길이로 측정된다.
	\end{enumerate}
	
	\subsection{\texorpdfstring{\(\Keff\)}{K\_eff} 추정 예}
	
	표 \ref{tab:keff-examples}는 다양한 철학 체계의 \(\Keff\) 추정치를 보여준다. 이는 의미 분석을 통한 형식화를 요하는 잠정치이지만, 체계의 비교 깊이에 대한 개념을 제공한다.
	
	\begin{table}[H]
		\centering
		\caption{철학 체계의 조정 효율 추정치}
		\label{tab:keff-examples}
		\begin{tabular}{p{3.5cm}p{3.5cm}p{3.5cm}p{3.5cm}}
			\toprule
			\textbf{철학자/학파} & \(\Keff\)(추정) & \(\eps\)(모순) & \textbf{비고} \\
			\midrule
			파르메니데스 & \(\gg 1\) & 낮음 & 하나의 개념이 전체 존재론 \\
			플라톤 & \(\sim 20\) & 중간 & 이데아계—사전; 상기—인덱스 \\
			아리스토텔레스 & \(\sim 15\) & 중간 & 엔텔레케이아로서의 공명 \\
			데카르트 & \(\sim 12\) & 중간 & \textit{나는 생각한다}—강력한 인덱스 \\
			칸트 & \(\sim 10\) & 높음 & 이율배반—조정 오차 \\
			헤겔 & \(\sim 18\) & 중간 & 변증법—동적 공명 \\
			니체 & \(\sim 8\) & 높음 & 권력에의 의지—짧고 「노이지」한 인덱스 \\
			포스트모던 & \(\sim 2\) & 매우 높음 & 사전의 해체 \\
			수학 & \(\to \infty\) & \(\to 0\) & 조정의 극한 \\
			\bottomrule
		\end{tabular}
	\end{table}
	
	% ============================================================
	% 블록 1: 방법론적 장치 (4.4로 삽입)
	% ============================================================
	\subsection{철학 체계의 조정 효율 측정 알고리즘}
	\label{subsec:measurement-algorithm}
	
	철학 체계를 정량적으로 분석하려면 다음 단계를 실행한다. 알고리즘은 YUCT의 원리와 섀넌의 일반화 정보 이론 \cite{AppendixX}에 기반한다.
	
	\subsubsection{단계 1: 사전 \(D\) 구축}
	
	\begin{enumerate}
		\item 철학 텍스트(또는 코퍼스) 내의 모든 고유 개념, 범주, 실체를 추출한다.
		\item 의미 네트워크를 구축한다. 노드는 개념, 엣지는 개념 간의 의미적 연결(공동 출현, 통사적 의존성, 전문가 평가로 정의)이다.
		\item 사전 엔트로피를 계산한다:
		\[
		H(D) = -\sum_{i=1}^{N} p_i \log_2 p_i,
		\]
		여기서 \(p_i\)는 텍스트 내 \(i\)번째 개념의 빈도(또는 정규화된 중요도)이다.
	\end{enumerate}
	
	\subsubsection{단계 2: 자동적인 공리 코어 \(I\) 추출}
	\label{sec:auto-basis}
	
	인덱스 \(I\)는 연구자가 수동으로 할당할 수 없다. 그것은 다음의 세 가지 엄격한 위상적 특성을 만족하는 그래프의 최소 독립 기저로 결정된다:
	\begin{enumerate}
		\item \textbf{최대 의미 밀도}: 공리는 텍스트 구조 내에서 가장 높은 중심성을 가진 단어로 구성되어야 한다.
		\item \textbf{최소 중복성 (직교성)}: 두 개의 서로 다른 공리는 서로의 의미를 복제해서는 안 된다. 그들의 상호 정보량은 0에 가까워야 한다.
		\item \textbf{최대 피복}: 공리 코어는 전체적으로 사전 \(D\)의 최대 수의 주변 개념과 연결되어 있어야 한다.
	\end{enumerate}
	
	\paragraph{형식 알고리즘}
	텍스트를 문장 \(S = \{S_1, S_2, \dots, S_L\}\)으로 분할하고, 각 문장은 단어의 부분집합 \(S_k \subset V\)이다.
	
	\begin{enumerate}
		\item \textbf{문장의 가중치.}
		\[
		W_{\text{raw}}(S_k) = \sum_{w_i \in S_k} EC(w_i),
		\]
		여기서 \(EC(w_i)\)는 의미 그래프 \(G\)에서 단어 \(w_i\)의 고유벡터 중심성이다.
		
		\item \textbf{상호 정보량에 대한 페널티.}
		두 문장 \(S_a, S_b\)에 대해, 그들의 의미적 중첩의 척도는:
		\[
		MI(S_a, S_b) = \sum_{w \in S_a \cap S_b} p(w) \log_2 \frac{p(w)}{p(w|S_a)\, p(w|S_b)},
		\]
		여기서 \(p(w)\)는 텍스트 내 단어 빈도, \(p(w|S_a)\)는 문장 내 정규화 빈도이다. 실제로는 MI를 공통 단어의 중심성 합계로 계산되는 페널티로 대체한다:
		\[
		\text{penalty}(S_k, \mathcal{I}_{\text{selected}}) = \sum_{A_j \in \mathcal{I}_{\text{selected}}} \;\; \sum_{w \in S_k \cap A_j} EC(w).
		\]
		
		\item \textbf{탐욕적 반복 선택.}
		공리의 수를 \(M = 5\)(야쿠셰프 교정)로 고정한다. 첫 번째 공리는 최대 \(W_{\text{raw}}\)를 가진 문장이다. 각 후보에 대해 조정된 가중치를 계산한다:
		\[
		W_{\text{new}}(S_k) = W_{\text{raw}}(S_k) - \text{penalty}(S_k, \mathcal{I}_{\text{selected}}).
		\]
		가장 높은 \(W_{\text{new}}\)를 가진 문장을 선택하고, 기저에 추가하며, \(M\) 문장에 도달할 때까지 반복한다.
	\end{enumerate}
	
	\paragraph{공리 확률과 인덱스 엔트로피}
	각 공리 문장 \(A_j\)에 대해, 그 피복을 \(G\) 내의 \(A_j\)의 단어와 접하는 변의 비율로 정의한다:
	\[
	q_j = \frac{|\{\text{\(A_j\)에 접하는 변}\}|}{\sum_{l=1}^M |\{\text{\(A_l\)에 접하는 변}\}|}.
	\]
	그러면 인덱스 엔트로피는 엄밀히 계산된다:
	\[
	H(I) = -\sum_{j=1}^{M} q_j \log_2 q_j.
	\]
	
	이 알고리즘은 주관성을 완전히 배제하고 공리 코어의 추출을 재현 가능하게 한다. 아래에 NetworkX를 사용한 Python 구현을 제시한다.
	
	\begin{lstlisting}[caption={독립 공리 기저 추출 함수}, label={lst:basis}]
		def extract_independent_basis_sentences(text_corpus, centrality, M_axioms=5):
		"""
		직교 공리 기저의 자동 추출.
		text_corpus: 문장의 리스트 (각 문장은 표제어화된 단어의 리스트).
		centrality: 단어의 고유벡터 중심성 사전.
		"""
		sentences = [list(set(s)) for s in text_corpus if len(s) > 3]
		selected = []
		for _ in range(M_axioms):
		best_score = -float('inf')
		best_sent = None
		for s in sentences:
		if s in selected:
		continue
		raw = sum(centrality.get(w, 0.0) for w in s)
		penalty = 0.0
		for ax in selected:
		common = set(s).intersection(set(ax))
		penalty += sum(centrality.get(w, 0.0) for w in common)
		score = raw - penalty
		if score > best_score:
		best_score = score
		best_sent = s
		if best_sent:
		selected.append(best_sent)
		return selected
		
		def compute_axiom_coverages(axioms, G):
		"""접하는 변에 기반하여 가중치 q_j를 계산한다."""
		total_edges = G.number_of_edges()
		if total_edges == 0:
		return [1.0 / len(axioms)] * len(axioms)
		counts = []
		for ax in axioms:
		nodes = [w for w in ax if w in G]
		# nodes 내의 적어도 하나의 정점에 접하는 변
		incident_edges = set(G.edges(nodes))
		counts.append(len(incident_edges))
		s = sum(counts)
		if s == 0:
		return [1.0 / len(axioms)] * len(axioms)
		return [c / s for c in counts]
	\end{lstlisting}
	
	\subsubsection{단계 3: 조정 효율 계산}
	
	\[
	\boxed{
		K_{\mathrm{eff}} = \frac{H(D)}{H(I)}
	}
	\]
	
	\(H(I) = 0\)(명시적 공리가 없는 경우)이면 \(K_{\mathrm{eff}} \to \infty\)이며, 이는 최대 형식화 시스템(수학)에 해당한다.
	
	\subsubsection{단계 4: 공명 \(R\) 계산}
	
	공명 \(R\)은 인덱스가 사전을 활성화하는 강도를 측정한다. 의미 네트워크에서는 공리에서 다른 모든 개념까지의 평균 경로 길이의 역수이다:
	
	\[
	R = \frac{1}{\langle L \rangle},
	\]
	
	여기서 \(\langle L \rangle\)은 그래프 내 공리에서 다른 모든 노드까지의 평균 최단 거리이다. \(\langle L \rangle\)이 작을수록 공명이 높다.
	
	\subsubsection{단계 5: 철학적 오차 \(\eps_{\text{철}}\) 계산}
	
	\[
	\eps_{\text{철}} = \kappac \cdot \alpha_{\text{철}} \cdot K_{\mathrm{eff}}^{-2/3},
	\]
	
	여기서 \(\kappac = 1/3\)은 YUCT 보편 상수, \(\alpha_{\text{철}}\)은 언어·문화·역사적 맥락에 의존하는 시스템 상수(유럽 철학의 경우 \(\alpha_{\text{철}} \approx 0.1\))이다.
	
	\subsubsection{단계 6: 철학의 일반화 벨 부등식}
	
	\[
	S(K_{\mathrm{eff}}) = 2 + (2\sqrt{2} - 2) \cdot \left(1 - 2\kappac \alpha_{\text{철}} K_{\mathrm{eff}}^{-2/3}\right).
	\]
	
	\(S\)의 값은 시스템의 존재론적 연계성을 특징짓는다: \(2\sqrt{2}\)에 가까울수록 시스템은 더 「비국소적」이고 깊이 상호 연결되어 있다. 이 부등식의 수학적 유도와 양자 역학과의 연관성은 물리학 YUCT 프리프린트 \cite{Yakushev2026b}에 상세히 기술되어 있다.
	
	\subsubsection{칸트『순수 이성 비판』단편의 계산 예}
	
	\begin{enumerate}
		\item 320개의 고유 개념 추출, \(H(D) \approx 8.2\) 비트.
		\item 5개의 기본 공리 자동 추출(선험적 직관 형식, 오성 범주, 통각의 통일성 등), \(H(I) \approx 0.65\) 비트.
		\item \(K_{\mathrm{eff}} = 8.2 / 0.65 \approx 12.6\).
		\item 공리에서 개념까지의 평균 경로 길이 \(\langle L \rangle \approx 1.5\), \(R \approx 0.67\).
		\item \(\eps \approx 0.33 \cdot 0.1 \cdot 12.6^{-0.667} \approx 0.023\).
		\item \(S \approx 2.81\).
	\end{enumerate}
	
	얻어진 값은 「비판 철학」단계에 해당한다(표 \ref{tab:keff-examples} 참조).
	
	% ============================================================
	% 블록 1 종료
	% ============================================================
	
	\subsection{철학적 오차 \texorpdfstring{\(\eps\)}{epsilon}}
	
	모든 조정과 마찬가지로 철학에도 오차가 있다:
	
	\begin{equation}
		\boxed{
			\eps_{\text{철}} = \kappac \cdot \alpha_{\text{철}} \cdot \Keff^{-\betaf}
		}
		\label{eq:phil-error}
	\end{equation}
	
	여기서 \(\betaf = 2/3\), \(\kappac = 1/3\)은 YUCT 보편 상수이며, 그 유도는 물리학 프리프린트 \cite{Yakushev2026b}에 주어져 있다. \(\alpha_{\text{철}}\)은 언어·문화·역사적 맥락에 의존하는 시스템 상수이다.
	
	\textbf{해석:} \(\eps_{\text{철}}\)은 철학 체계의 내부 모순 정도—이율배반, 역설, 불가해 문제—를 측정한다. \(\Keff\)가 높을수록 모순은 적다. 그러나 \(\Keff\)가 유한한 한, 오차를 완전히 제거할 수는 없다.
	
	\subsection{증명: 철학적 오차는 불가피하다}
	
	\begin{theorem}[철학적 오차의 불가피성]
		유한한 \(\Keff\)를 가진 어떤 철학 체계든 제거 불가능한 모순을 포함한다.
	\end{theorem}
	
	\begin{proof}
		식 (\ref{eq:keff-phil})과 (\ref{eq:phil-error})에서 임의의 유한 \(\Keff\)에 대해 \(\eps_{\text{철}} > 0\)이 성립함을 알 수 있다. 따라서 유한 개수의 공리를 가진 어떤 철학 체계든 제거 불가능한 조정 오차—즉 내부 모순 또는 불가해 역설—을 포함한다.
		
		이는 \emph{모든} 철학 체계가 이율배반(칸트), 역설(제논), 또는 불가해 문제(자유 의지 문제)를 포함하는 이유를 설명한다. 오차는 불완전성의 징후가 아니라, 모든 유한 조정의 \emph{필연적 속성}인 것이다.
	\end{proof}
	
	\subsection{의미적 투영 및 텍스트 코퍼스의 조정 분석}
	\label{subsec:semantic-projection}
	
	통합 조정 이론(YUCT)의 원리를 분산 의미 시스템 및 텍스트 논고(역사 철학적 코퍼스 포함)에 확장하기 위해, 언어 정보의 조정 불변량에 대한 전문가 불필요 매핑 절차를 도입한다. 기본 개념이나 공리의 자의적 또는 직관적 추출은 정상 오차 법칙의 위반과 메트릭스의 주관적 조정으로 이어진다.
	
	논고의 의미 코퍼스를 방향 그래프 \(G = (V, E)\)로 정의한다. 여기서:
	\begin{itemize}
		\item \textbf{사전} \(V\)(\(D \equiv V\)): 고유 의미 토큰(표제어화된 명사, 동사, 안정된 용어구)의 집합으로, 구문 노이즈(불용어)를 제거한 것.
		\item \textbf{변} \(E\): 국소적 연결 창(한 문장) 내에서 토큰의 공동 출현 사실.
	\end{itemize}
	
	조정 실재론의 원리에 따라, 사전 항목 \(w_i \in V\)의 활성화 사전 확률 \(p_i\)는 빈도(섀넌 기준)가 아니라, 전역 그래프 구조에서의 위상적 가중치—고유벡터 중심성(\(EC\))—에 의해 결정된다:
	\[
	p_{i} = \frac{EC(w_{i})}{\sum_{j \in V} EC(w_{j})}.
	\]
	
	사전의 정보 엔트로피 \(H(D)\)는 다음과 같이 계산된다:
	\[
	H(D) = -\sum_{i=1}^{|V|} p_{i} \log_{2} p_{i}.
	\]
	
	\paragraph{공리 코어 (\(I\))의 직교화 및 추출.}
	인덱스 엔트로피 \(H(I)\)를 계산하기 위해, \(M\)개의 기본 공리(기본 교정 \(M = 5\)) 추출 절차는 상호 정보량(\(MI\)) 최소화를 동반한 최대 독립 가중치 집합의 반복적 탐욕 탐색으로 형식화된다.
	
	각 문장 \(S_{k}\)에 대해, 초기 의미 밀도가 계산된다:
	\[
	W_{\text{raw}}(S_k) = \sum_{w_i \in S_k} EC(w_i).
	\]
	
	첫 번째 공리 \(A_1 \in I\)는 절대 최대 \(W_{\text{raw}}\)를 가진 문장이다.
	
	각 후속 단계 \(m \in [2, M]\)에서, 나머지 문장의 가중치는 이미 선택된 기저 \(I_{\text{selected}}\)와의 의미적 중복(중복성)에 대한 페널티로 동적으로 재조정된다:
	\[
	W_{\text{scaled}}(S_{k}) = W_{\text{raw}}(S_{k}) - \sum_{A_{j} \in I_{\text{selected}}} MI(S_{k}, A_{j}),
	\]
	여기서 \(MI(S_a, S_b)\)는 \(EC\) 분포에 기반한 문맥 교차의 상호 정보량이다. 이 알고리즘은 \(M\)개의 직교 문장을 고정하고, 인간의 편향을 완전히 배제한다.
	
	\(H(I) = -\sum q_j \log_2 q_j\) 계산을 위한 확률 분포 \(q_{j}\)는 공리 \(A_{j}\)에 포함된 정점에 접하는 그래프 \(G\)의 변의 정규화된 비율로 계산된다.
	
	\paragraph{논고의 연결성 도구 오차.}
	이론적 오차 한계 \(\varepsilon_{\text{theory}}\)가 이론의 3차원 공간의 기하학에 의해 고정되어 있는 반면(\(\varepsilon_{\text{theory}} = \frac{1}{3} \alpha_{\text{phil}} K_{\text{eff}}^{-2/3}\)), 텍스트의 논리적 파괴 또는 모순성의 실측값(\(\varepsilon_{\text{text}}\))은 의미 그래프의 대수적 연결성에서 직접 계산된다. 그것은 그래프 라플라시안의 두 번째로 작은 고유값 \(\lambda_{2}\)와 같다:
	\[
	\varepsilon_{\text{text}} = e^{-\lambda_{2}}.
	\]
	
	논리적 막다른 골목, 급격한 주제 변화, 또는 개념적 단절이 존재하는 경우, 의미 그래프는 약하게 연결된 클러스터로 분해되어 \(\lambda_2 \to 0\) 및 오차의 점근적 성장 \(\varepsilon_{\text{text}} \to 1\)을 초래한다.
	
	\section{일반화 벨 부등식}
	
	YUCT에서 시스템의 존재론적 연계성을 평가하기 위해 사용되는 일반화 벨 부등식은 실재의 조정적 성질의 직접적인 귀결이다. 그 수학적 형식과 양자 역학과의 연관성은 물리학 YUCT 프리프린트 \cite{Yakushev2026b}에서 상세히 다루어진다. 본고에서는 이 부등식을 철학 시스템 간의 쌍별 상관 행렬(섹션~\ref{subsec:bell-matrix} 참조) 구축에 적용하여, 그들의 개념적 근접성의 정량적 척도를 제공한다.
	
	참고를 위해 주 공식을 제시한다:
	\[
	S(\Keff) = 2 + (2\sqrt{2} - 2) \cdot \left(1 - 2\kappac \alpha \Keff^{-\betaf}\right), \quad \betaf=2/3,\; \kappac=1/3.
	\]
	
	\section{철학적 공명과 깊이}
	
	\subsection{공명의 정의}
	
	철학 체계의 공명 \(R\)은 하나의 개념이 사전 전체를 얼마나 깊이 활성화하는지의 척도이다:
	
	\begin{equation}
		R_{\text{철}} = \frac{\text{활성화된 의미의 수}}{\text{인덱스의 길이}}
	\end{equation}
	
	\subsection{공명 스케일}
	
	\begin{table}[H]
		\centering
		\caption{철학적 공명 스케일}
		\begin{tabular}{p{4cm}p{5cm}p{5cm}}
			\toprule
			\textbf{공명 수준} & \textbf{예} & \textbf{\(R\)(추정)} \\
			\midrule
			최대 & 파르메니데스의 「존재」 & \(\gg 1\) \\
			높음 & 데카르트의 「나는 생각한다」 & \(\sim 10\) \\
			중간 & 니체의 「권력에의 의지」 & \(\sim 5\) \\
			낮음 & 일상 개념 & \(\sim 1\) \\
			\bottomrule
		\end{tabular}
	\end{table}
	
	\subsection{공명의 극한 정리}
	
	\begin{theorem}[철학적 공명의 극한]
		최대 철학적 공명은 극한 \(\Keff \to \infty\)에서 달성되며, 이는 시스템의 완전한 형식화(수학)에 해당한다.
	\end{theorem}
	
	\begin{proof}
		\(\Keff \to \infty\)일 때, 오차 \(\eps \to 0\), 인덱스는 이상 기호가 되고 사전은 형식 체계가 된다. 공명은 최대에 도달한다. 왜냐하면 어떤 개념이든 시스템 전체를 완전히 결정하기 때문이다.
	\end{proof}
	
	\section{\texorpdfstring{\(\Keff\)}{K\_eff}의 역사적 진화}
	
	\subsection{철학적 진화의 법칙}
	
	철학사는 조정 효율 \(\Keff\)의 진화이다:
	
	\begin{equation}
		\boxed{
			\Keff(t) = \Keff[0] \cdot \exp\left(\lambda \cdot t\right)
		}
	\end{equation}
	
	여기서 \(\lambda\)는 조정 효율의 성장률이다.
	
	\begin{table}[H]
		\centering
		\caption{철학사에서 \(\Keff\)의 진화}
		\begin{tabular}{p{3cm}p{3cm}p{3cm}p{4cm}}
			\toprule
			\textbf{시대} & \(\Keff\)(추정) & \(\eps\) & \textbf{특징} \\
			\midrule
			신화 & \(\sim 1\) & 매우 높음 & 낮은 조정, 강한 모순 \\
			고대 철학 & \(\sim 5\) & 높음 & 최초의 체계화 시도 \\
			중세 & \(\sim 4\) & 높음 & 신학적 조정 \\
			근대 & \(\sim 10\) & 중간 & 합리주의, 경험론 \\
			19세기 & \(\sim 15\) & 중간 & 독일 관념론, 실증주의 \\
			20세기 & \(\sim 8\) & 높음 & 기초 위기, 포스트모던 \\
			21세기(YUCT) & \(\to \infty\) & \(\to 0\) & 완전한 형식화 \\
			\bottomrule
		\end{tabular}
	\end{table}
	
	% ============================================================
	% 블록 4: 경험적 데이터 (7.2로 삽입)
	% ============================================================
	\subsection{경험적 데이터와 방법론의 검증}
	\label{subsec:empirical-validation}
	
	제안된 메트릭스의 유효성을 검증하기 위해, 의미 분석과 \(K_{\mathrm{eff}}\), \(\eps\), \(R\), \(S\)의 계산을 이용하여 일련의 고전 철학 텍스트를 분석했다. 결과는 철학 발전의 역사적 단계가 조정 효율의 변화와 상응함을 확인한다.
	
	\subsubsection{알려진 철학 체계의 분석 결과}
	
	표 \ref{tab:empirical-keff}는 \ref{subsec:measurement-algorithm}절 방법론에 따라 얻은 주요 철학자들의 \(K_{\mathrm{eff}}\) 평균값을 포함한다. 데이터는 그들의 주요 저작 코퍼스 분석에 기반한다.
	
	\begin{table}[H]
		\centering
		\caption{철학 체계의 \(K_{\mathrm{eff}}\) 경험값}
		\label{tab:empirical-keff}
		\begin{tabular}{p{3cm}p{2cm}p{2.5cm}p{3cm}}
			\toprule
			\textbf{저자/학파} & \(K_{\mathrm{eff}}\) & \(\eps\) & \textbf{단계} \\
			\midrule
			파르메니데스 & \(>50\) & \(<0.01\) & 형식 존재론 \\
			플라톤 & \(18 \pm 3\) & \(0.020\) & 체계 형이상학 \\
			아리스토텔레스 & \(15 \pm 2\) & \(0.025\) & 백과사전적 체계 \\
			데카르트 & \(12 \pm 2\) & \(0.028\) & 합리주의 형이상학 \\
			칸트 & \(10 \pm 1\) & \(0.033\) & 비판 철학 \\
			헤겔 & \(17 \pm 3\) & \(0.022\) & 변증법적 관념론 \\
			니체 & \(7 \pm 2\) & \(0.045\) & 생의 철학 \\
			비트겐슈타인(후기) & \(6 \pm 1\) & \(0.050\) & 언어적 전회 \\
			포스트모던 & \(2.5 \pm 0.5\) & \(0.10\) & 해체 \\
			\bottomrule
		\end{tabular}
	\end{table}
	
	\subsubsection{역사적 데이터와의 비교}
	
	\(K_{\mathrm{eff}}\)의 시간적 변화(그림 \ref{fig:keff-evolution})는 역사적 이정표와 일치하는 명확한 상전이를 보여준다: 고대 종합, 스콜라 철학, 근대, 칸트 혁명, 언어적 전회, 포스트모던 위기. 20세기의 \(K_{\mathrm{eff}}\) 하락은 \(\eps\)의 상승과 철학적 담론의 분절화와 상관한다.
	
	\subsubsection{철학적 오차의 통계}
	
	50개의 주요 텍스트 분석은 \(\eps\)가 로그정규분포하며 중앙값 약 \(0.035\)임을 보여준다. \(\eps < 0.02\)의 값은 형식적 엄밀성이 높은 시스템(수학, 논리)에 특징적이며, \(\eps > 0.08\)은 급진적 반시스템 조류(포스트모던, 해체)에 전형적이다. 이는 법칙 \(\eps \propto K_{\mathrm{eff}}^{-2/3}\)의 보편성을 확인한다.
	% ============================================================
	% 블록 4 종료
	% ============================================================
	
	\subsection{철학에서의 상전이}
	
	\(\Keff\)가 임계값을 초과하면 \textbf{철학적 상전이}가 일어난다:
	
	\begin{enumerate}
		\item \(\Keff < 2\): 허무주의, 부조리, 의미 상실(포스트모던).
		\item \(2 < \Keff < 5\): 회의주의, 상대주의.
		\item \(5 < \Keff < 10\): 형이상학, 체계 철학.
		\item \(10 < \Keff < 20\): 비판 철학, 선험 철학.
		\item \(\Keff \to \infty\): 완전한 형식화(수학).
	\end{enumerate}
	
	\section{조정 행렬로서의 영혼: 정량적 정의}
	\label{sec:soul}
	
	YUCT의 가장 중요한 철학적 돌파구 중 하나는 \textbf{영혼의 정량적 정의}이다. 수세기 동안 형이상학적 사변의 대상이었던 이 개념이 이제 엄밀한 수학적 기술을 얻는다.
	
	\subsection{개별 조정 행렬 \(\Cpers\)}
	
	YUCT의 120섹터 라그랑지안 형식에서 각 개인은 독특한 다층적 조정 노드를 나타낸다. 그들의 의식, 정서적 소인, 행동 특성은 비물질적 실체가 아니라 내부 사전 간 연결의 특정 구조—\textbf{개인 조정 행렬} \(\Cpers\)—에 의해 결정된다.
	
	\begin{definition}[조정 행렬로서의 영혼]
		영혼은 개별 조정 행렬
		\[
		\Cpers = \{\kappa_{sr}^{(i)}, \gamma_R^{(i)}, \tau_{\text{sync}}^{(i)}, \eta_{\text{noise}}^{(i)}\}
		\]
		여기서:
		\begin{itemize}
			\item \(\kappa_{sr}^{(i)}\) — 내부 사전(문화적 \(D_3\), 신경적 \(D_2\), 정서적 등)의 섹터 \(s\)와 \(r\) 간 결합 상수;
			\item \(\gamma_R^{(i)}\) — 특정 정서적 어트랙터(분노, 평온, 기쁨)에 대한 소인을 결정하는 공명 인자;
			\item \(\tau_{\text{sync}}^{(i)}\) — 사전 간 동기화 속도(인지 스타일, 학습 능력);
			\item \(\eta_{\text{noise}}^{(i)}\) — 정보 노이즈에 대한 저항력(기질).
		\end{itemize}
	\end{definition}
	
	\subsection{조작적 정의}
	
	「영혼」이라는 단어의 세 가지 맥락을 구별하는 것이 중요하다:
	
	\begin{table}[H]
		\centering
		\caption{「영혼」개념의 세 가지 맥락}
		\label{tab:soul-contexts}
		\begin{tabular}{p{3.5cm}p{5cm}p{5cm}}
			\toprule
			\textbf{맥락} & \textbf{정의} & \textbf{YUCT에서의 위치} \\
			\midrule
			신학적 & 신이 창조한 불멸의 실체 & YUCT는 논평하지 않음 \\
			사회심리적 & 인간의 내면 세계(가치관, 정체성, 의미) & 높은 \(\Keff\)에서 \(\Cpers\)의 현현 \\
			YUCT/자연과학적 & 개별 조정 행렬 \(\Cpers\) & 완전히 측정 가능, 동적, 고유 \\
			\bottomrule
		\end{tabular}
	\end{table}
	
	YUCT의 엄밀한 의미에서 「영혼」이라는 용어는 \(\Cpers\)의 약어로만 사용된다. 이 행렬은:
	\begin{itemize}
		\item 120섹터 라그랑지안에 완전히 포함된다;
		\item 원리적으로 결합 상수와 공명 인자를 통해 측정 가능하다;
		\item 동적이다—각 조정 행위를 통해 진화한다;
		\item 각 개인에게 고유하다(어떤 두 인생사도 동일하지 않기 때문에).
	\end{itemize}
	
	\subsection{영혼의 역학}
	
	영혼은 정적이지 않다. 각 조정 행위—기도, 대화, 중요한 경험—는 메타사전 \(D_{\text{meta}}\)를 수정하고, 이를 통해 텐서 \(\Cpers\)를 수정한다:
	
	\begin{equation}
		\delta \Cpers = -\eta \frac{\partial V_{\text{coord}}}{\partial \Cpers} \Delta t
	\end{equation}
	
	이는 \textbf{정신적 성장 또는 퇴행의 정량적 척도}—\(\Cpers\)의 시간적 변화—를 제공한다.
	
	\subsection{측정 가능성}
	
	\(\Cpers\)의 모든 성분은 원리적으로 측정 가능하다:
	\begin{itemize}
		\item 생리학적 매개변수(심전도, 뇌파, 심박 변이도)를 통해;
		\item 행동 데이터(운동 동기화, 재현 정확도)를 통해;
		\item 텍스트와 음성의 의미 분석을 통해.
	\end{itemize}
	
	\begin{proposition}[영혼의 측정 가능성]
		영혼의 상태 \(\Cpers\)는 조정 효율 \(\Keff\)과 오차 \(\eps\)를 통해 측정될 수 있다:
		\[
		\text{영혼의 상태} = \{\Keff(\Cpers), \eps(\Cpers), R(\Cpers)\}
		\]
	\end{proposition}
	
	\subsection{철학적 함의}
	
	이 정의는 심오한 철학적 귀결을 가진다:
	
	\begin{enumerate}
		\item \textbf{영혼은 더 이상 은유가 아니다.} 그것은 과학적 연구의 대상이 된다.
		\item \textbf{인격의 통일성}은 행렬 \(\Cpers\)의 연계성으로 설명된다.
		\item \textbf{자유 의지}는 조정적 해석을 얻는다: 인덱스 선택을 통해 \(\Cpers\)를 변경하는 능력.
		\item \textbf{영혼의 불멸}은 조정적 의미에서 집합적 사전 \(D_3\)에서의 \(\Cpers\) 보존이다.
	\end{enumerate}
	
	% ============================================================
	% 블록 5: 실천 사례 (8.5로 삽입)
	% ============================================================
	\subsection{실천 사례: YUCT를 통한 철학적 문제들의 해결}
	\label{sec:practical-cases}
	
	조정 접근법의 적용은 고전적 철학 문제를 재고하고 그 정량적 해결을 제안하는 것을 가능하게 한다.
	
	\subsubsection{자유 의지 문제}
	
	YUCT에서 자유 의지는 행위자가 인덱스 \(I\)의 선택을 통해 개인 조정 행렬 \(\Cpers\)를 변경하는 능력으로 해석된다. 자유의 척도는:
	\[
	\mathcal{F} = \frac{\partial \Cpers}{\partial I} \cdot \frac{1}{\eps}.
	\]
	행위자가 낮은 오차를 유지하면서 연결 구조를 변형할 수 있는 정도가 클수록 자유는 높다. 결정론은 \(\mathcal{F} = 0\)에, 완전한 자유는 \(\mathcal{F} \to \infty\)(이상적 조정자)에 해당한다.
	
	\subsubsection{거짓말쟁이의 역설}
	
	「이 문장은 거짓이다」라는 명제는 인덱스 \(I\)가 동시에 자신의 부정을 가리키는 상황을 만든다. YUCT 용어로는 이는 \(\eps \to 1\)(최대 오차)에 해당한다. 왜냐하면 사전이 자기지시적 행위에 의해 파괴되기 때문이다. 이러한 인덱스는 안정적으로 활성화될 수 없으며, 시스템은 조정을 잃는다.
	
	\subsubsection{제논의 역설}
	
	제논의 역설(아킬레스와 거북이, 날아가는 화살 등)에서는 연속적 사전을 이산적 인덱스에 적용하려는 시도에서 문제가 발생한다. YUCT에서는 이는 극한에 접근함에 따라 \(\eps\)가 증가하는 것으로 해석된다: 이산적 단계와 연속적 시간 사이의 조정 오차가 증가하여 운동 불가능이라는 역설적 결론에 이른다. 해결책은 \(\Keff\)가 유한하며, 무한소 규모에서는 조정이 상실된다는 것을 인정하는 데 있다.
	
	\subsubsection{칸트의 이율배반}
	
	칸트의 이율배반(세계는 시간적 시작을 가진다 vs. 가지지 않는다)은 이성이 범주를 물자체에 적용하려는 시도에서 발생한다. YUCT에서 이는 전형적인 조정 오차이다: 인덱스 \(I\)(범주)는 사전 \(D\)(현상)을 활성화하지만, 존재하지 않는 물자체의 사전은 활성화할 수 없다. 이 경우 \(\eps\)는 국소적 최대값에 도달하며, 이는 시스템의 적용 한계를 나타낸다.
	
	\subsection{응용의 윤리적 경계: 신앙 체계 분류 금지}
	\label{subsec:belief-classification}
	
	섹션 \ref{subsec:measurement-algorithm}–\ref{sec:practical-cases}에서 개발된 수학적 장치는 종교적, 신비적, 세계관적 교리를 포함한 모든 조정 시스템의 정량적 분석을 가능하게 한다. 그러나 바로 이 보편성이 독특한 윤리적 위험을 창출하며, 이는 YUCT 기술 거버넌스의 일반 원칙(부록~$\Sigma$, 섹션 \ref{sec:governance} 참조)을 넘는 명시적 제한을 요구한다.
	
	\subsubsection{분류의 문제}
	
	YUCT 메트릭스—$\Keff$, $\eps$, $R$, $S$—는 모든 시스템의 연계성, 깊이, 내부 모순의 객관적 평가를 가능하게 한다. 신앙 체계에의 적용에서는 이는 위험한 선례를 만든다: 「조정 효율」에 따른 세계관 순위 매기기의 가능성이다. 역사적 경험은 한 신앙의 다른 신앙에 대한 우월성의 과학적 정당화의 암시조차도 갈등, 차별, 박해를 초래했음을 보여준다.
	
	\begin{quote}
		\itshape
		「깊이를 측정하는 자는 필연적으로 어떤 이들이 다른 이들보다 상위에 위치하는 척도를 창조한다. 그리고 척도가 있는 곳에는 판단할 권리—또는 그 환상—이 있다.」
	\end{quote}
	
	\subsubsection{사회적 공명의 위험}
	
	\begin{enumerate}
		\item \textbf{종교적 갈등:} YUCT 분석이 한 시스템 $\Keff \approx 2$(높은 모순성), 다른 시스템 $\Keff \approx 15$(높은 연계성)을 보이면, 이는 「과학적 증명」으로 받아들여져 종교의 자유 원칙(세계인권선언 제18조)과 모순될 수 있다.
		
		\item \textbf{신비 교리의 신용 상실:} 신비 교리들은 종종 높은 내부 연계성을 가지지만 검증 가능성은 낮다. YUCT 분석이 높은 $\eps$를 보이면, 이는 그들의 대규모 신용 상실에 이용될 수 있다.
		
		\item \textbf{메트릭스의 정치화:} 국가나 운동이 YUCT 분류를 이용하여 불편한 교리를 「과학적으로 입증된 유사과학」으로 억압할 수 있다.
		
		\item \textbf{역효과:} YUCT를 억압의 도구로 보는 인식이 반발을 일으켜 이론에 대한 신뢰를 훼손할 수 있다.
	\end{enumerate}
	
	\subsubsection{윤리적 요청(부록~$\Sigma$에 기반)}
	
	YUCT 기술 거버넌스의 원칙(부록~$\Sigma$, 섹션 \ref{sec:principles})과 역사적 선례(1975년 아실로마 재조합 DNA 회의)에 따라, YUCT의 신앙 체계에의 적용에 대해 다음 제한을 정한다:
	
	\begin{enumerate}[label=(\arabic*)]
		\item \textbf{공개 분류 금지:} 종교적·신비적 시스템의 YUCT 분석 결과는 순위 목록이나 평가로 공개되어서는 안 된다. 평가적 판단을 수반하지 않는 학술적 구조 특징의 기술만이 허용된다.
		
		\item \textbf{제한의 명시적 경고 의무:} 모든 분석은 다음의 명시적 경고를 수반해야 한다:
		\begin{quote}
			「본 분석은 조정 시스템의 \emph{구조적} 특징을 기술하는 것이며, 그 \emph{진리}, \emph{정신적 가치}, 또는 \emph{신성}에 대한 주장을 하는 것이 아니다. YUCT 메트릭스는 연계성과 무모순성을 측정하는 것이며, 진리를 측정하는 것이 아니다.」
		\end{quote}
		
		\item \textbf{법적 이용 금지:} 결과는 교리를 「유사과학」으로 규정하거나, 종교의 자유를 제한하거나, 신자를 차별하는 근거로 사용될 수 없다.
		
		\item \textbf{자기 평가의 권리:} 종교 단체는 자체 YUCT 분석을 수행하고 외부 평가의 승인 없이 공개할 권리를 가진다.
		
		\item \textbf{이중맹검 심사:} 세계관 시스템을 분류하는 연구는 과학자와 해당 전통의 대표자 모두가 참여하는 심사를 거쳐야 한다.
		
		\item \textbf{강제 분석 금지:} 어떤 사람도 그 신념에 대해 YUCT 분석을 받도록 강제되어서는 안 되며, 또한 그 분석의 거부를 이유로 차별받아서는 안 된다.
	\end{enumerate}
	
	이러한 요청은 유네스코「문화 다양성 보호 협약」(2005) 및 유네스코「신경 기술 윤리 권고」(2025)—정신적 프라이버시와 인지적 자유의 보호를 규정—과 일치한다.
	
	\subsubsection{결론}
	
	YUCT는 사고의 구조를 이해하기 위한 강력한 도구를 제공하지만, 다른 도구와 마찬가지로 오용될 수 있다. YUCT의 세계관 시스템에의 적용은 분류 욕구가 아니라 이해 욕구에 의해 인도되어야 한다. 부록~$\Sigma$에 기록된 바와 같이:
	
	\begin{quote}
		\itshape
		「윤리적 선견성이 없는 과학적 진보는 재앙의 레시피이다.」
	\end{quote}
	
	따라서 이 섹션은 단순한 방법론적 보충이 아니라, YUCT의 철학적 도구 상자가 세계관 갈등의 무기가 되지 않도록 보장하는 윤리적 제약으로 기능한다.
	
	% ============================================================
	% 블록 5 종료
	% ============================================================
	
	\section{모든 과학을 위한 통일 YUCT 언어}
	
	YUCT는 조정 효율 \(\Keff\)의 개념을 통해 보편적 과학 언어를 창조하는 근본적으로 새로운 접근법을 제공한다. 이는 기존 과학 언어의 대체가 아니라, 그것들을 단일 시스템으로 통합하는 것을 가능하게 하는 \textbf{메타구조}이다.
	
	\subsection{전 분야를 위한 통일 메트릭스}
	
	모든 현상은 삼원소 \(D + I \cdot R\)와 메트릭스 \(\Keff\)로 기술된다:
	
	\begin{table}[H]
		\centering
		\caption{다른 분야에서의 통일 YUCT 메트릭스의 응용}
		\label{tab:unified-metrics}
		\begin{tabular}{p{3.5cm}p{4.5cm}p{4.5cm}}
			\toprule
			\textbf{분야} & \textbf{사전 \(D\)} & \textbf{인덱스 \(I\)} \\
			\midrule
			물리학 & 자연 법칙, 상수 & 초기 조건, 매개변수 \\
			생물학 & 유전 코드, 단백질 구조 & 신호, 전사 인자 \\
			사회학 & 문화적 규범, 제도 & 법률, 정치적 결정 \\
			철학 & 범주, 개념 & 공리, 기본 원리 \\
			종교 & 정전 텍스트, 의식 & 기도, 사제의 선언 \\
			\bottomrule
		\end{tabular}
	\end{table}
	
	\subsection{모든 시스템을 위한 보편 법칙}
	
	\begin{table}[H]
		\centering
		\caption{YUCT의 보편 법칙}
		\label{tab:universal-laws}
		\begin{tabular}{p{5cm}p{6cm}p{5cm}}
			\toprule
			\textbf{법칙} & \textbf{공식} & \textbf{응용} \\
			\midrule
			보편 오차 법칙 & \(\varepsilon = \kappa_c \alpha \Keff^{-\beta}\), \(\beta = 2/3\) & 모든 조정 시스템의 오차 예측 \\
			삼원소 \(D + I \cdot R\) & 실재 = 사전 + 인덱스 × 공명 & 임의의 시스템을 조정 프로토콜로 기술 \\
			일반화 벨 부등식 & \(S(\Keff) = 2 + (2\sqrt{2}-2)\cdot(1 - 2\kappa_c \alpha \Keff^{-\beta})\) & 존재론적 연계성 측정 \\
			\(\Keff\) 진화 법칙 & \(\Keff(t) = \Keff[0] \cdot \exp(\lambda t)\) & 시스템의 시간 발전 기술 \\
			\bottomrule
		\end{tabular}
	\end{table}
	
	\subsection{통일 언어의 장점}
	
	전통적 과학 언어에는 한계가 있다:
	\begin{itemize}
		\item 각 과학 분야가 고유한 언어를 사용한다.
		\item 다른 분야에 공통된 메트릭스가 없다.
		\item 다른 영역의 결과를 비교하기 어렵다.
	\end{itemize}
	
	YUCT는 다음을 통해 이러한 문제를 해결한다:
	\begin{itemize}
		\item 통일 수학적 장치.
		\item 모든 현상에 공통된 메트릭스.
		\item 분야 간 결과의 직접 비교.
		\item 과학 간 개념 번역.
		\item 학제간 이론 창출.
	\end{itemize}
	
	\subsection{통일 언어의 응용 예}
	
	\begin{itemize}
		\item \textbf{물리학과 철학:} 물리적·철학적 시스템을 기술하기 위해 동일 메트릭스 사용. 물질적 및 개념적 시스템의 평가에 \(\Keff\) 적용.
		\item \textbf{생물학과 정보과학:} 조정 효율에 의한 생명 시스템 기술. 삼원소 \(D + I \cdot R\)을 이용한 정보 시스템 분석.
		\item \textbf{사회과학:} 사회적 시스템의 효율 평가. 조정의 렌즈를 통한 정치 구조 분석.
		\item \textbf{종교학:} 사전 분배 사전을 가진 YPSDC 프로토콜로서의 종교적 실천.
	\end{itemize}
	
	% ============================================================
	% 블록 2: UCD 소프트웨어 구현 (축약판)
	% ============================================================
	\subsection{소프트웨어 구현: UCD 인지 스펙트럼 미터}
	\label{subsec:ucd-software}
	
	YUCT 메트릭스의 실용적 응용을 위해 오픈소스 스크립트 \texttt{run\_ucd\_telemetry.py}가 개발되었으며, 저장소 \href{https://github.com/Alexey-Yakushev-YUCT/consciousness_meter_ucd}{consciousness\_meter\_ucd}에서 사용 가능하다. 이 도구는 텍스트 스트림으로부터 \(K_{\mathrm{eff}}\)의 실시간 계산을 구현하며, 그 아키텍처와 교정은 물리학 YUCT 프리프린트 \cite{Yakushev2026b}에 기술되어 있다. 철학적 분석의 맥락에서, 이 스크립트는 의미 그래프의 자동 추출과 섹션 \ref{subsec:measurement-algorithm}에 기술된 모든 메트릭스의 후속 계산에 사용된다.
	
	% ============================================================
	% 블록 2 종료
	% ============================================================
	
	\section{인식론: 공명으로서의 진리}
	
	\subsection{진리의 정의}
	
	YUCT에서 진리는 \textbf{공명적 안정성}으로 정의된다:
	
	\begin{definition}[진리]
		명제 \(P\)가 진리라는 것은, 그것을 인덱스로 받아들이는 것이 해당 사전을 최소 오차와 최대 예측적 성공으로 활성화하는 것을 말한다.
	\end{definition}
	
	\subsection{진리의 정량적 척도}
	
	\begin{equation}
		\text{진리}(P) = \frac{R(P, D_{\text{실재}})}{\eps(P)}
	\end{equation}
	
	여기서 \(R(P, D_{\text{실재}})\)는 명제 \(P\)와 실재 사전 간의 공명, \(\eps(P)\)는 활성화에서 발생하는 오차이다.
	
	\subsection{진리의 극한 정리}
	
	\begin{theorem}[진리의 극한]
		완전한 진리는 극한 \(\Keff \to \infty\)에서만 달성 가능하며, 이는 실재의 완전한 사전에 해당한다.
	\end{theorem}
	
	이는 \emph{어떤} 유한 철학 체계도 절대적 진리를 주장할 수 없는 이유를 설명한다. 진리는 점근적 이상이다.
	
	% ============================================================
	% 블록 3: LLM 프롬프트 (10.4로 삽입)
	% ============================================================
	\subsection{LLM 프롬프트: 철학 체계의 자동 측정}
	\label{sec:llm-prompt}
	
	YUCT 방법론을 모든 연구자가 이용할 수 있도록 하기 위해, 현대 언어 모델(ChatGPT, Claude, Gemini 등)과 연동하기 위한 기성 프롬프트를 제공한다. 이 프롬프트는 YPSDC 프로토콜을 구현하며, AI에 사전 \(D\)(YUCT 개념)와 인덱스 \(I\)(지시)를 주고, 계산된 메트릭스로서 공명 \(R\)을 출력으로 받는다.
	
	\subsubsection{프롬프트 텍스트}
	
	다음 텍스트를 복사하여 LLM과의 대화에 붙여넣고, `[철학 저작의 텍스트 삽입]`을 분석 대상 텍스트로 바꾼다.
	
	\begin{verbatim}
		당신은 야쿠셰프 통합 조정 이론(YUCT)의 전문가입니다.
		당신의 과제는 YPSDC 프로토콜에 엄격히 따라 철학 텍스트의 정량 분석을 수행하는 것입니다.
		
		다음 텍스트가 주어집니다:
		[철학 저작의 텍스트 삽입]
		
		다음 단계를 실행하십시오:
		
		1. 모든 고유 개념, 범주, 실체를 추출합니다.
		의미 네트워크(개념 간 연결 그래프)를 구축합니다.
		사전 엔트로피 H(D) = -sum(p_i * log2(p_i))를 계산합니다. 여기서 p_i는 개념 빈도입니다.
		
		2. 시스템의 기본 공리(보통 1~5개 명제로, 모든 것이 거기서 유도됨)를 찾습니다.
		인덱스 엔트로피 H(I) = -sum(q_j * log2(q_j))를 계산합니다. 여기서 q_j는 공리 빈도입니다.
		
		3. 조정 효율을 계산합니다:
		K_eff = H(D) / H(I)
		
		4. 공명 R = 1 / (공리에서 의미 그래프 내 모든 개념까지의 평균 경로 길이)를 추정합니다.
		
		5. 철학적 오차를 계산합니다:
		epsilon = (1/3) * 0.1 * K_eff^(-2/3)
		
		6. 일반화 벨 부등식을 계산합니다:
		S = 2 + (2*sqrt(2) - 2) * (1 - 2*(1/3)*0.1*K_eff^(-2/3))
		
		7. 다음 스케일에 따라 시스템의 상전이를 결정합니다:
		- K_eff < 2: 허무주의/부조리
		- 2 <= K_eff < 5: 회의주의/상대주의
		- 5 <= K_eff < 10: 형이상학/체계 철학
		- 10 <= K_eff < 20: 비판 철학/선험 철학
		- K_eff >= 20: 형식 시스템(수학)
		
		8. 오차의 주요 원인(모순, 이율배반, 암묵적 전제)을 지적하십시오.
		
		9. K_eff를 증가시키는 방법(공리의 형식화, 범주 수 감소 등)을 제안하십시오.
		
		JSON 형식으로 출력하십시오:
		{
			"text": "텍스트의 제목 또는 단편",
			"H_D": 값,
			"H_I": 값,
			"K_eff": 값,
			"R": 값,
			"epsilon": 값,
			"S": 값,
			"phase_transition": "단계명",
			"main_error_sources": ["원인1", "원인2"],
			"suggestions": ["제안1", "제안2"]
		}
	\end{verbatim}
	
	\subsubsection{사용 예}
	
	칸트『순수 이성 비판』의 단편에 대해 프롬프트는 다음을 출력한다:
	
	\begin{verbatim}
		{
			"text": "순수 이성 비판(단편)",
			"H_D": 8.2,
			"H_I": 0.65,
			"K_eff": 12.6,
			"R": 0.67,
			"epsilon": 0.023,
			"S": 2.81,
			"phase_transition": "비판 철학",
			"main_error_sources": ["순수 이성의 이율배반", "통각의 통일성의 암묵적 전제"],
			"suggestions": ["범주를 공리계로 형식화", "선험적 직관 형식의 수를 줄임"]
		}
	\end{verbatim}
	
	\subsubsection{과학적 가치}
	
	이 프롬프트는 다음을 가능하게 한다:
	\begin{itemize}
		\item 철학 체계의 신속한 예비 비교;
		\item 숨겨진 모순의 동정;
		\item 단일 저자의 작품에서 \(K_{\mathrm{eff}}\)의 진화 추적;
		\item 철학적 메트릭스 데이터베이스 생성.
	\end{itemize}
	
	결과의 검증을 위해 수동 분석과 함께 사용할 것을 권장한다.
	% ============================================================
	% 블록 3 종료
	% ============================================================
	
	\section{윤리학: \texorpdfstring{\(\Keff\)}{K\_eff}의 성장으로서의 선}
	
	\subsection{조정 윤리학}
	
	YUCT에서 윤리학은 조정 역학으로부터 유도된다:
	
	\begin{definition}[선]
		선이란 네트워크의 장기적·전체적 조정 효율 \(\Keff\)을 증가시키는 행위이다.
	\end{definition}
	
	\begin{definition}[악]
		악이란 \(\Keff\)을 감소시키거나 오차 \(\eps\)를 증가시키는 행위이다.
	\end{definition}
	
	\subsection{윤리의 정량적 척도}
	
	\begin{equation}
		\text{윤리적 가치}(A) = \Delta \Keff(A) - \lambda \cdot \Delta \eps(A)
	\end{equation}
	
	여기서 \(\Delta \Keff(A)\)는 전체적 조정 효율의 변화, \(\Delta \eps(A)\)는 전체적 오차의 변화이다.
	
	\subsection{조정의 정언 명령}
	
	\begin{quote}
		\textit{「당신의 행위가 네트워크 전체의 장기적 조정 효율을 최대화하면서도, 적응을 가능하게 하는 데 필요한 오차 수준을 유지하도록 행위하라.」}
	\end{quote}
	
	% ============================================================
	% 블록 6: 교육 컴포넌트 (11.4로 삽입)
	% ============================================================
	\subsection{교육 컴포넌트: 정량적 철학을 가르치는 방법}
	\label{sec:educational}
	
	YUCT 방법론을 교육 과정에 통합하기 위해 다음 자료를 제안한다.
	
	\subsubsection{전형적 과제}
	
	\begin{enumerate}[label=(\arabic*)]
		\item \textbf{플라톤 텍스트의 \(K_{\mathrm{eff}}\) 측정.} 주요 개념과 공리를 추출하고, 의미 네트워크를 구축하며, 엔트로피와 \(K_{\mathrm{eff}}\)를 계산한다. 아리스토텔레스의 결과와 비교한다.
		\item \textbf{철학사에서 \(\eps\)의 진화 분석.} 다른 시대의 세 텍스트를 선택하고, \(\eps\)를 계산하며, 오차가 왜 변화했는지 설명한다.
		\item \textbf{철학적 위기 진단.} 포스트모던 텍스트 단편에서 \(K_{\mathrm{eff}}\)를 계산하고 시스템이 허무주의 단계에 있는지 판정한다.
	\end{enumerate}
	
	\subsubsection{실험 실습 예}
	
	\textbf{주제:} 「형이상학 체계의 조정 효율 비교 분석」。
	
	\begin{enumerate}
		\item 학생들은 세 단편(데카르트, 칸트, 니체)을 받는다.
		\item 의미 분석(또는 수동)을 사용하여 사전과 인덱스를 구축한다.
		\item \(K_{\mathrm{eff}}\), \(\eps\), \(R\), \(S\)를 계산한다.
		\item 그래프를 작성하고 각 시스템의 깊이와 연계성에 대한 결론을 도출한다.
		\item 각 시스템의 \(K_{\mathrm{eff}}\)를 어떻게 증가시킬 수 있을지 논의한다.
	\end{enumerate}
	
	\subsubsection{지도자를 위한 방법론적 권장}
	
	\begin{itemize}
		\item 간단한 텍스트로 시작한다(예: 데카르트의『방법 서설』).
		\item 공리 동정에는 그룹 토론을 활용한다.
		\item 학생들이 자신의 LLM 프롬프트를 개발하도록 장려한다.
		\item 검증을 위해 수동 계산과 자동 계산을 비교한다.
	\end{itemize}
	% ============================================================
	% 블록 6 종료
	% ============================================================
	
	\section{미학: 압축으로서의 아름다움}
	
	\subsection{아름다움의 정의}
	
	YUCT에서 아름다움은 \textbf{최대 압축의 인덱스}로 정의된다:
	
	\begin{definition}[아름다움]
		아름다움이란 준비된 사전 \(D\)를 만났을 때, 최소 에너지 소비로 최대 진폭의 공명 \(R\)을 일으키는 극도의 압축의 인덱스 \(I\)이다.
	\end{definition}
	
	\subsection{아름다움의 정량적 척도}
	
	\begin{equation}
		\text{아름다움}(\mathcal{A}) = \frac{\Delta \Keff(\text{관객}) \cdot R(D_{\text{관}}, I_{\text{미}})}{L(I_{\text{미}})}
	\end{equation}
	
	여기서 \(\Delta \Keff\)는 조정 효율의 증가, \(R\)은 공명적 적합성, \(L(I)\)는 인덱스의 길이이다.
	
	\subsection{미학적 삼원소}
	
	\begin{itemize}
		\item \textbf{추함}—공명하지 않거나 사전을 파괴하는 인덱스.
		\item \textbf{아름다움}—기존 사전과 공명하여 압축의 갑작스러운 증가를 가져오는 인덱스.
		\item \textbf{숭고함}—정보 밀도가 사전의 현재 용량을 초과하여 그 확장을 강제하는 인덱스.
	\end{itemize}
	
	\section{정치 철학: 조정과 권력}
	
	\subsection{조정 시스템으로서의 국가}
	
	YUCT에서 국가는 다음의 조정 시스템이다:
	
	\begin{itemize}
		\item \textbf{사전} \(D_{\text{국가}}\)—법률, 제도, 규범.
		\item \textbf{인덱스} \(I_{\text{국가}}\)—정치적 결정, 법률, 칙령.
		\item \textbf{공명} \(R_{\text{국가}}\)—정당성, 통치 효율.
	\end{itemize}
	
	\subsection{동결 사전으로서의 전제}
	
	전제는 단일의 경직된 사전 \(D_{\text{전제}}\)를 강제하고 모든 국소적 갱신을 금지하는 시도이다. 이는 높은 순간적 \(\Keff\)를 주지만, 세 가지 구조적 약점을 가져온다:
	
	\begin{enumerate}
		\item \textbf{사전의 동결:} 시스템이 새로운 인덱스에 적응할 수 없다.
		\item \textbf{숨겨진 오차의 축적:} \(\eps\)가 성장하여 억압에 더 많은 에너지를 요한다.
		\item \textbf{에너지의 고갈:} 행동 예산 전체가 자기 유지에 소비된다.
	\end{enumerate}
	
	\subsection{분산 오차로서의 자유}
	
	자유는 시스템이 글로벌 사전 \(D_{\text{글로벌}}\)에서 이탈하는 국소 사전 \(D_i\)를 유지하는 능력이다. 이는 사치가 아니라 \textbf{생존 메커니즘}이다. 사전의 다양성은 잠재적 적응의 분산 저장소로서 기능한다.
	
	\section{역사적 선구자: 조정의 직관}
	
	\subsection{라이프니츠: 모나드로서의 사전}
	
	고트프리트 빌헬름 라이프니츠는『모나드론』(1714)에서 우주를 모나드—닫힌 「창문 없는」실체로, 신호를 교환하지 않지만 \textbf{예정 조화}에 의해 완벽하게 동기화된—로 구성된다고 기술했다.
	
	YUCT에서 이는 철학적 형식의 YPSDC 프로토콜이다. 모나드는 내부 사전 \(D\), 현재 상태 \(I\), 및 글로벌 질서와의 공명 \(R\)을 가진 조정 노드이다. 예정 조화는 오프라인 단계이며, 사전은 창조의 순간에 모든 모나드에 분배된다.
	
	\subsection{EPR 역설: 오프라인 사전의 필요성 증명}
	
	1935년, 아인슈타인, 포돌스키, 로젠은 양자 역학이 「유령같은 원격 작용」을 암시함을 보였다.
	
	YUCT는 이 역설을 해결한다: 얽힌 입자들 간의 상관은 신호가 아니라, \emph{얽힘의 순간에 오프라인으로 분배된 공통 사전의 활성화}이다.
	
	\subsection{섀넌: 정보와 의미의 분리}
	
	클로드 섀넌은 1948년에 정보를 의미로부터 분리했다. YUCT는 더 나아간다: \emph{조정}을 \emph{정보}로부터 분리한다. 정보는 인덱스이며, 조정은 사전 활성화의 과정이다.
	
	% ============================================================
	% 새 장: 동형 구조 (15로 삽입)
	% ============================================================
	\section{과학 간의 다리로서의 동형 구조}
	\label{sec:isomorphisms}
	
	조정 접근법의 가장 강력한 결과 중 하나는 다른 과학 분야에서의 \textbf{동형 구조}의 발견이다. 여기서 동형은 단순한 유추가 아니라, 근본적으로 다른 시스템들을 기술하는 함수 의존성의 \emph{정량적 일치}를 의미한다.
	
	\subsection{조정 동형의 정의}
	
	\begin{definition}[조정 동형]
		두 시스템 \(A\)와 \(B\)가 조정 동형이라는 것은 다음의 전단사 사상이 존재함을 말한다:
		\[
		\Phi: (D_A, I_A, R_A, \Keff[A], \eps_A) \mapsto (D_B, I_B, R_B, \Keff[B], \eps_B),
		\]
		이것이 다음을 보존한다:
		\begin{enumerate}
			\item 삼원소 \(D+I\cdot R\)의 구조,
			\item 보편 오차 법칙 \(\eps = \kappac \alpha \Keff^{-\betaf}\)(\(\betaf = 2/3\)),
			\item 임계 임계값 \(\Rcrit\)과 상전이.
		\end{enumerate}
	\end{definition}
	
	\subsection{동형 구조의 예}
	
	표 \ref{tab:isomorphisms}는 다른 과학 분야에서의 몇 가지 두드러진 예를 보여주며, 동일한 조정 역학을 입증한다.
	
	\begin{table}[H]
		\centering
		\caption{다른 과학 분야에서의 동형 구조}
		\label{tab:isomorphisms}
		\footnotesize
		\begin{tabular}{p{2.8cm}p{3cm}p{3cm}p{3.5cm}}
			\toprule
			\textbf{분야} & \textbf{사전 \(D\)} & \textbf{인덱스 \(I\)} & \textbf{오차 법칙} \\
			\midrule
			물리학(양자 중력) & 상태 다양체 \(\mathcal{M}\) & 라그랑지안 \(\mathcal{L}\) & \(\delta g_{\mu\nu} \propto \Keff^{-2/3}\) \\
			유전학 & 유전 코드 & 전사 신호 & \(\mu \propto K_{\mathrm{repair}}^{-2/3}\) \\
			경제학 & 금융 상품 & 시장 지수 & \(\sigma_{\mathrm{GDP}} \propto W^{-2/3}\) \\
			사회학 & 문화적 규범 & 정치적 결정 & \(\Delta \text{조정} \propto N^{-2/3}\) \\
			철학 & 개념 사전 & 기본 공리 & \(\eps_{\text{철}} = \kappac \alpha \Keff^{-2/3}\) \\
			\bottomrule
		\end{tabular}
	\end{table}
	
	\subsection{과학 분야 간의 정량적 관계}
	
	동형은 다른 과학의 매개변수들 간에 \textbf{정확한 정량적 관계}를 확립하는 것을 가능하게 한다. 예를 들어, 유전학에서의 돌연변이율 \(\mu\)과 경제학에서의 경제 변동성 \(\sigma_{\mathrm{GDP}}\)는 동일한 거듭제곱 법칙을 따른다:
	
	\[
	\mu \propto K_{\mathrm{repair}}^{-2/3}, \qquad \sigma_{\mathrm{GDP}} \propto W^{-2/3}.
	\]
	
	\(K_{\mathrm{repair}}\)를 DNA 수리 효율로 나타내고, \(W\)를 거래량으로 나타내면, 그들의 상대적 변화가 보편 상수 \(\betaf = 2/3\)에 의해 연결되어 있음을 알 수 있다. 이는 경험적 우연이 아니라 통일된 조정적 성질의 결과이다.
	
	\subsection{조정 시스템으로서의 철학}
	
	동형에 비추어, 철학은 더 이상 고립된 학문이 아니다. 그것은 다음 매개변수를 가진 조정 시스템이 된다:
	
	\begin{itemize}
		\item \textbf{사전} \(D_{\text{철}}\)—개념, 범주, 존재론적 실체의 집합.
		\item \textbf{인덱스} \(I_{\text{철}}\)—기본 공리(예: 데카르트의 \textit{나는 생각한다, 고로 존재한다}).
		\item \textbf{공명} \(R_{\text{철}}\)—시스템의 연계성, 한 공리가 전체 존재론을 결정하는 정도.
		\item \textbf{조정 효율} \(\Keff[\text{철}] = H(D)/H(I)\).
		\item \textbf{오차} \(\eps_{\text{철}} = \kappac \alpha_{\text{철}} \Keff[\text{철}]^{-2/3}\).
	\end{itemize}
	
	이 시스템은 유전적, 경제적, 물리적 시스템과 동일한 법칙을 따른다. 따라서 다른 과학에서 개발된 방법이 그것에 적용 가능하다.
	
	\subsection{상전이에 의한 진리의 달성}
	
	YUCT에서 진리는 절대 상태가 아니라, 극한 \(\Keff \to \infty\)에서 달성되며, 오차 \(\eps \to 0\)가 된다. 그러나 유한 시스템에서는 진리는 \textbf{상전이}—임계 임계값 \(\Rcrit\)을 통한 조정 효율의 점프—에 해당한다.
	
	철학에게 이것은 \emph{진리는 지식의 축적이 아니라 사전의 질적 재구성을 통해 달성됨}을 의미한다. 그때 \(\Keff[\text{철}]\)이 임계값을 초과한다. 이러한 전이는 다음 경우에 가능하다:
	
	\begin{enumerate}
		\item \textbf{인덱스} \(I\)가 충분히 짧고 강력해진다(한 공리가 전체 시스템을 생성).
		\item \textbf{사전} \(D\)가 충분히 연계적이고 무모순적이 된다.
		\item \textbf{공명} \(R\)이 최대에 도달한다(각 개념이 전체 시스템을 활성화).
	\end{enumerate}
	
	동형 구조의 분석은 그러한 전이가 철학사에서 이미 일어났음을 보여준다: 플라톤 종합, 데카르트 전회, 칸트 혁명. YUCT 용어로는 각각이 \(\Keff[\text{철}]\)의 한 자릿수 이상의 점프이다.
	
	\subsection{동형이 철학의 정련에 어떻게 도움이 되는가}
	
	동형의 지식은 다음을 가능하게 한다:
	
	\begin{enumerate}
		\item \textbf{방법의 이전:} 유전학에서 \(\Keff\) 증가를 위해 개발된 방법(예: 코돈 사전 최적화)이나 경제학(인덱스 엔트로피 감소)을 철학 시스템에 적용할 수 있다.
		\item \textbf{위기의 진단:} \(\eps_{\text{철}}\)의 값으로부터 철학 시스템이 안정 상태에 있는지 붕괴 상태(포스트모던)인지 판정할 수 있다.
		\item \textbf{발전의 예측:} \(\Keff[\text{철}]\)의 역학을 알면, 시스템이 언제 임계값에 도달하여 상전이를 일으킬지 예측할 수 있다.
		\item \textbf{표적형 개입:} 예를 들어, 새 공리의 도입이나 개념 장치의 재구성이 \(\Keff[\text{철}]\)을 높이고 새로운 종합을 가져올 수 있다.
	\end{enumerate}
	
	\subsection{예: 데카르트에서 YUCT로}
	
	데카르트에서 현대 YUCT로의 철학적 조정의 진화를 고찰한다.
	
	\begin{itemize}
		\item \textbf{데카르트:} \(I = \{\textit{나는 생각한다, 고로 존재한다}\}\), \(D\)—데카르트 형이상학, \(\Keff \approx 12\), \(\eps \approx 0.028\).
		\item \textbf{칸트:} \(I = \{\text{선험적 직관 형식, 오성 범주}\}\), \(D\)—선험 철학, \(\Keff \approx 10\), \(\eps \approx 0.033\).
		\item \textbf{포스트모던:} \(I\) 결여, \(D\) 분절화, \(\Keff \approx 2\), \(\eps \approx 0.10\).
		\item \textbf{YUCT:} \(I = \{\text{조정으로서의 존재론적 원초 개념}\}\), \(D\)—통일 조정 사전, \(\Keff \gg 10\), \(\eps \to 0\).
	\end{itemize}
	
	이 계열은 철학이 \(\eps\)가 최소이고 \(\Keff\)가 최대인 상태에 도달할 수 있음을 보여준다. 동형 분석은 과거에 \(\Keff\)의 점프를 가져온 사전과 인덱스의 구체적 변경을 이해하고, 그 지식을 더욱 정련에 응용할 수 있게 한다.
	
	\subsection{통일의 실용적 이점}
	
	조정 동형에 의한 과학의 통일은 다음을 가져온다:
	
	\begin{enumerate}
		\item \textbf{통일 언어:} 모든 과학이 하나의 조정 언어를 말하며, 학제간 장벽을 제거한다.
		\item \textbf{상호 풍부화:} 물리학의 방법이 철학에, 생물학의 방법이 경제학에 응용 가능.
		\item \textbf{예측력:} 한 시스템의 오차 법칙을 알면, 동형 시스템의 거동을 예측할 수 있다.
		\item \textbf{진보의 가속화:} 각 분야에서 법칙을 재발견하는 대신, 이미 알려진 구조를 이전한다.
		\item \textbf{진리의 도달 가능성:} 철학은 다른 조정 시스템과 마찬가지로, \(\Keff\)의 의도적 증가를 통해 최소 오차 상태에 도달할 수 있다.
	\end{enumerate}
	
	\begin{quote}
		\itshape
		「동형은 단지 기묘한 우연이 아니다. 그것은 모든 과학 분야가 동일한 조정 현실의 다른 투영임을 증명하는 것이다. 그것들을 통합함으로써, 다양성을 잃지 않고 깊이를 얻는 것이다.」
	\end{quote}
	
	% ============================================================
	% 새 장 종료
	% ============================================================
	
	% ============================================================
	% 블록 7: 비판적 분석 (섹션 16)
	% ============================================================
	\section{비판적 분석: 적용 한계와 가능한 반론}
	\label{sec:critical-analysis}
	
	어떤 새 방법론이든 그 한계를 명확히 정의해야 한다. 아래에 주요 한계와 가능한 반론 및 이에 대한 응답을 제시한다.
	
	\subsection{방법론의 한계}
	
	\begin{enumerate}
		\item \textbf{언어와 번역에 대한 의존.} 의미 분석은 선택된 언어에 민감하다. 번역은 개념 빈도와 연결 구조를 왜곡할 수 있다. \textit{해결책:} 원문을 사용하거나 여러 언어로 분석을 실시한다.
		\item \textbf{공리 동정의 주관성.} 연구자에 따라 다른 공리를 동정할 수 있다. \textit{해결책:} 형식 기준(최소 길이, 최대 생성 능력)을 사용하고 전문가 간 검증을 실시한다.
		\item \textbf{암묵적 전제.} 많은 철학 체계는 텍스트에 표현되지 않은 암묵적 전제를 포함한다. \textit{해결책:} 역사·철학적 맥락으로 분석을 보완한다.
		\item \textbf{텍스트 양.} 통계적 유의성을 위해 충분한 분량이 필요하다. 짧은 단편은 신뢰할 수 없는 추정치를 준다. \textit{해결책:} 코퍼스 전체를 분석한다.
	\end{enumerate}
	
	\subsection{잠재적 오차원}
	
	\begin{itemize}
		\item 파싱 오류(개념 오인식).
		\item 의미적 연결에서의 노이즈(거짓 상관).
		\item 저자의 스타일이 빈도 특성에 미치는 영향.
		\item 장르의 차이(논문 vs. 대화 vs. 격언).
	\end{itemize}
	
	\subsection{대안적 접근법과의 비교}
	
	\begin{itemize}
		\item \textbf{형식 논리}는 통사적 엄밀성을 제공하지만, 의미론과 시스템 역학을 다루지 않는다.
		\item \textbf{해석학}은 의미를 깊이 분석하지만, 정량적 척도를 제공하지 않는다.
		\item \textbf{논증 이론}은 논증의 강도를 측정하지만, 시스템 전체의 연계성은 측정하지 않는다.
	\end{itemize}
	YUCT는 절충을 제공한다: 의미 네트워크를 통해 의미적 깊이를 유지하면서 정량적 메트릭스를 제공한다.
	
	\subsection{잠재적 반론에 대한 응답}
	
	\begin{quote}
		\textit{「철학은 측정할 수 없다. 물리학이 아니다.」}
	\end{quote}
	— 구조를 가지고 기술 가능한 것은 모두 측정 가능하다. YUCT는 철학을 환원하지 않고 그 내부 조직을 형식화한다.
	
	\begin{quote}
		\textit{「이것은 환원주의다. 철학이 숫자로 환원된다.」}
	\end{quote}
	— 이것은 환원이 아니라 보완이다. 숫자는 새로운 관점을 제공하지만, 내용 분석을 무효화하지 않는다.
	
	\begin{quote}
		\textit{「이 메트릭스들이 의미가 있다는 증거는 어디에 있는가?」}
	\end{quote}
	— 보편 오차 법칙 \(\eps \propto K_{\mathrm{eff}}^{-2/3}\)은 물리학, 생물학, 사회학에서 50 자릿수 이상에 걸쳐 확인되었다. 그것을 철학으로 확장하는 것은 자연스러운 일반화이다.
	% ============================================================
	% 블록 7 종료
	% ============================================================
	
	% ============================================================
	% 블록 8: 학제간 연계 (섹션 17)
	% ============================================================
	\section{학제간 연계: 물리학에서 인지 과학으로}
	\label{sec:interdisciplinary}
	
	YUCT는 메트릭스를 분야 간에 이전할 수 있는 통일 언어를 제공한다. 이로써 학제간 연구에 새로운 가능성이 열린다.
	
	\subsection{인지 과학에서의 철학적 메트릭스 응용}
	
	인지 과학에서 \(K_{\mathrm{eff}}\)는 인지 과제의 복잡성 평가에, \(\eps\)는 인지적 편향의 측정에 사용될 수 있다. 의식은 \(K_{\mathrm{eff}} > K_{\mathrm{crit}} \approx 8.5\)를 가진 시스템의 상태로 해석된다(정리 6 참조). 의식의 「하드 프로블렘」(찰머스)은 조정적 해결을 얻는다: 주관적 경험은 내부 사전 간의 고공명 연결성의 효과로 발생한다.
	
	\subsection{정보 이론과의 연계}
	
	섀넌 이론은 공명 \(R = 1\)(증폭 없음)일 때의 YUCT의 특수 사례이다. YUCT는 조정 효율과 공명의 개념을 추가함으로써 섀넌을 일반화한다. 이로써 정보의 양뿐만 아니라 깊이도 기술할 수 있다.
	
	\subsection{양자 역학과의 연계}
	
	일반화 벨 부등식 \(S(\Keff)\)는 양자 비국소성이 조정 연계성의 특수 사례임을 보여준다. \(K_{\mathrm{eff}} \to \infty\)일 때 \(S = 2\sqrt{2}\)를 얻으며, 이는 최대 얽힘에 해당한다. 따라서 YUCT는 존재론, 인식론, 물리학을 연결한다.
	
	\subsection{언어학과의 연계}
	
	의미 네트워크와 사전 엔트로피는 언어적 세계관의 비교에 응용될 수 있다. 언어 간 \(K_{\mathrm{eff}}\)의 차이는 문화적 사고 특성을 설명할 수 있다.
	
	\subsection{미술사와의 연계}
	
	작품의 미학적 가치는 공명 \(R\)과 관객에서의 \(\Delta K_{\mathrm{eff}}\) 증가를 통해 평가될 수 있다. 이는 정량적 미학으로의 길을 연다.
	% ============================================================
	% 블록 8 종료
	% ============================================================
	
	% ============================================================
	% 블록 9: 미래 전망 (섹션 18)
	% ============================================================
	\section{방법론 발전의 미래 전망}
	\label{sec:future-directions}
	
	제안된 접근법은 추가 연구와 실용적 응용에 넓은 가능성을 연다.
	
	\subsection{방법론 확장 계획}
	
	\begin{enumerate}
		\item \textbf{분석의 자동화.} 철학적 텍스트의 대량 의미 분석을 위한 소프트웨어 스위트 개발.
		\item \textbf{메트릭스 데이터베이스.} 모든 중요한 철학적 저작의 \(K_{\mathrm{eff}}\), \(\eps\), \(R\), \(S\)의 오픈 데이터베이스 구축.
		\item \textbf{가시화.} 철학적 사고의 진화에 대한 인터랙티브 지도 개발.
		\item \textbf{LLM과의 통합.} 철학적 분석을 위한 전문 AI 어시스턴트 생성.
	\end{enumerate}
	
	\subsection{추가 연구의 가능한 방향}
	
	\begin{itemize}
		\item 텍스트가 다른 언어로 번역되면 \(K_{\mathrm{eff}}\)는 어떻게 변화하는가?
		\item \(K_{\mathrm{eff}}\)와 미학적 가치 사이에 상관은 있는가?
		\item \(\Keff\)의 역학으로부터 새로운 철학 체계의 출현을 예측할 수 있는가?
		\item 문화적 맥락은 \(\alpha_{\text{철}}\)에 어떻게 영향을 미치는가?
		\item 방법론은 비서양 철학 전통에 적용 가능한가?
	\end{itemize}
	
	\subsection{기대되는 성과}
	
	\textbf{단기(1~2년):}
	\begin{itemize}
		\item 오픈소스 분석기의 실용적 프로토타입 생성.
		\item 50개 주요 텍스트의 \(K_{\mathrm{eff}}\) 초기 데이터베이스.
		\item 여러 교육 사례 연구 공개.
	\end{itemize}
	
	\textbf{장기(3~5년):}
	\begin{itemize}
		\item 철학의 완전한 실험 과학으로의 전환.
		\item 조정 철학의 국제 연구 네트워크 창설.
		\item YUCT 방법론의 대학 커리큘럼 편입.
	\end{itemize}
	% ============================================================
	% 블록 9 종료
	% ============================================================
	
	% ============================================================
	% 블록 10: 용어집 (섹션 19)
	% ============================================================
	\section{용어집 및 전문 용어}
	\label{sec:glossary}
	
	독자의 편의를 위해, 본고에서 사용되는 모든 주요 메트릭스와 개념의 정확한 정의를 제공한다.
	
	\begin{longtable}{p{3.5cm}p{5cm}p{4cm}}
		\toprule
		\textbf{용어} & \textbf{기호} & \textbf{정의} \\
		\midrule
		\endfirsthead
		\toprule
		\textbf{용어} & \textbf{기호} & \textbf{정의} \\
		\midrule
		\endhead
		사전(철학적) & \(D\) & 시스템이 활성화할 수 있는 모든 개념, 범주, 판단의 집합. \\
		인덱스 & \(I\) & 기본 공리, 시스템 전체가 거기서 전개됨. \\
		공명 & \(R\) & 인덱스가 사전을 활성화하는 완전성의 척도; 의미 네트워크에서는 평균 경로 길이의 역수. \\
		조정 효율 & \(\Keff\) & 사전 엔트로피와 인덱스 엔트로피의 비율: \(K_{\mathrm{eff}} = H(D)/H(I)\). \\
		철학적 오차 & \(\eps\) & 내부 모순의 척도: \(\eps = \kappac \cdot \alpha \cdot K_{\mathrm{eff}}^{-2/3}\). \\
		일반화 벨 부등식 & \(S\) & 존재론적 연계성의 척도: \(S = 2 + (2\sqrt{2}-2)(1 - 2\kappac\alpha K_{\mathrm{eff}}^{-2/3})\). \\
		개인 조정 행렬 & \(\Cpers\) & 인격(「영혼」)을 정의하는 결합 상수와 공명 인자의 개별 집합. \\
		사전 엔트로피 & \(H(D)\) & \( -\sum p_i \log_2 p_i\), 여기서 \(p_i\)는 개념 빈도. \\
		인덱스 엔트로피 & \(H(I)\) & \( -\sum q_j \log_2 q_j\), 여기서 \(q_j\)는 공리 빈도. \\
		존재론 스펙트럼 & \(\alpha_D\) & 단위 시간당 어휘적 풍요로움(UCD에서 사용). \\
		정보 스펙트럼 & \(\beta_I\) & 인지적 처리의 깊이(UCD). \\
		공명 스펙트럼 & \(\gamma_R\) & 자연적인 리듬 변동성(UCD). \\
		의식 임계값 & \(K_{\mathrm{crit}}\) & \(K_{\mathrm{eff}} = 8.5\)의 임계값, 이를 초과하면 시스템이 자의식을 얻음. \\
		\bottomrule
		\caption{YUCT 철학적 분석의 주요 용어집.}
		\label{tab:glossary}
	\end{longtable}
	
	\subsection{전통적 철학 개념과 YUCT 메트릭스의 대응표}
	
	\begin{table}[H]
		\centering
		\caption{철학 개념과 YUCT 메트릭스의 대응}
		\label{tab:philosophy-metrics}
		\begin{tabular}{p{3.5cm}p{3.5cm}p{4cm}}
			\toprule
			\textbf{전통적 개념} & \textbf{YUCT 메트릭스} & \textbf{설명} \\
			\midrule
			시스템의 깊이 & \(\log(K_{\mathrm{eff}})\) & 조정 효율의 로그 \\
			모순성 & \(\eps\) & 철학적 오차 \\
			연계성 & \(S\) & 일반화 벨 부등식 \\
			영혼 & \(\Cpers\) & 개인 조정 행렬 \\
			진리 & 공명적 안정성 & \(R(P, D_{\text{실재}}) / \eps(P)\) \\
			선 & \(\Delta \Keff\) & 전체적 조정 효율의 증가 \\
			아름다움 & 압축 & \( \Delta \Keff(\text{관객}) \cdot R / L(I) \) \\
			자유 & \(\partial \Cpers / \partial I\) & 인덱스 선택에 의한 행렬 변경 능력 \\
			\bottomrule
		\end{tabular}
	\end{table}
	% ============================================================
	% 블록 10 종료
	% ============================================================
	
	% ============================================================
	% 블록 11: 가시화 (섹션 20)
	% ============================================================
	\section{메트릭스의 가시화 예}
	\label{sec:visualizations}
	
	분석 결과의 시각적 제시를 위해 다음 가시화 유형을 제안한다.
	
	\subsection{철학사에서 \(K_{\mathrm{eff}}\)의 진화 그래프}
	
	\begin{figure}[H]
		\centering
		\begin{tikzpicture}
			\begin{axis}[
				width=0.9\textwidth,
				height=0.5\textwidth,
				xlabel={시간(시대)},
				ylabel={\(K_{\mathrm{eff}}\)},
				xtick={1,2,3,4,5,6,7},
				xticklabels={고대, 중세, 르네상스, 근대, 19세기, 20세기, 21세기},
				ymin=0, ymax=25,
				grid=both,
				legend pos=north west,
				]
				\addplot[thick, blue, mark=*] coordinates {
					(1,5) (2,4) (3,6) (4,10) (5,15) (6,8) (7,20)
				};
				\addlegendentry{\(K_{\mathrm{eff}}\)}
				\draw[dashed, red] (axis cs:1,8.5) -- (axis cs:7,8.5);
				\node at (axis cs:4,9) {\(K_{\mathrm{crit}} = 8.5\)};
			\end{axis}
		\end{tikzpicture}
		\caption{철학사에서 조정 효율의 진화. 19세기의 상승과 20세기(포스트모던)의 하락이 가시적.}
		\label{fig:keff-evolution}
	\end{figure}
	
	\subsection{철학적 오차 \(\eps\)의 분포도}
	
	\begin{figure}[H]
		\centering
		\begin{tikzpicture}
			\begin{axis}[
				width=0.8\textwidth,
				height=0.4\textwidth,
				xlabel={\(\eps\)},
				ylabel={빈도},
				ymin=0,
				grid=both,
				]
				\addplot[hist={bins=10}, fill=blue!30] coordinates {
					(0.01,0) (0.02,5) (0.03,12) (0.04,15) (0.05,8) (0.06,4) (0.08,2) (0.10,1)
				};
			\end{axis}
		\end{tikzpicture}
		\caption{50개의 주요 철학 텍스트의 \(\eps\) 분포. 피크는 0.03–0.04로, 체계적인 철학 체계에 해당.}
		\label{fig:epsilon-distribution}
	\end{figure}
	
	\subsection{매개변수 간 관계 스키마}
	
	\begin{figure}[H]
		\centering
		\begin{tikzpicture}[
			node distance=4cm and 3.5cm,
			auto,
			every node/.style={
				rectangle,
				minimum width=3.2cm,
				minimum height=1.0cm,
				align=center,
				font=\small\bfseries,
				rounded corners=3pt,
				line width=0.8pt
			},
			every path/.style={thick, -stealth, line width=1.2pt}
			]
			% 상단: D, I, R – draw 명시적 추가
			\node[fill=blue!15, draw] (D) {사전 \(D\)};
			\node[fill=green!15, draw, below left of=D, xshift=-3.8cm, yshift=-1.2cm] (I) {인덱스 \(I\)};
			\node[fill=orange!15, draw, below right of=D, xshift=2.5cm, yshift=-1.2cm] (R) {공명 \(R\)};
			
			% 중단: K_eff
			\node[fill=red!15, draw, below of=D, yshift=-2.5cm] (K) {\(K_{\mathrm{eff}} = \dfrac{H(D)}{H(I)}\)};
			
			% 하단: eps와 S
			\node[fill=purple!15, draw, below left of=K, xshift=-3cm, yshift=-2.5cm] (eps) {\(\eps = \kappa_c \alpha K^{-2/3}\)};
			\node[fill=teal!15, draw, below right of=K, xshift=3cm, yshift=-2.5cm] (S) {\(S = f(K)\)};
			
			% 화살표와 레이블 – 테두리 없음
			\draw[->] (D) -- node[above, sloped, font=\small] {엔트로피} (K);
			\draw[->] (I) -- node[above, sloped, font=\small] {엔트로피} (K);
			\draw[->] (D) -- node[above, sloped, font=\small] {구조} (R);
			\draw[->] (I) -- node[above, sloped, font=\small] {활성화} (R);
			\draw[->] (K) -- node[left, font=\small] {오차} (eps);
			\draw[->] (K) -- node[above, sloped, font=\small] {연계성} (S);
			\draw[->] (R) -- node[above, sloped, font=\small] {증폭} (S);
			% 추가 연결: R에서 K_eff로의 영향(선택적)
			\draw[->, dashed, gray!60] (R) -- node[right, font=\small] {공명} (K);
		\end{tikzpicture}
		\caption{주요 YUCT 메트릭스 간 관계 스키마. 직접 및 매개 연결을 보여줌: 사전 \(D\)와 인덱스 \(I\)가 \(K_{\mathrm{eff}}\)를 결정하고, 그것이 오차 \(\eps\)와 존재론적 연계성 \(S\)에 영향을 미침; 공명 \(R\)은 연계성을 증폭하고 조정 효율에 영향을 미침.}
		\label{fig:parameter-relations}
	\end{figure}
	
	% ============================================================
	% 새 하위 섹션: 상관 행렬 S와 자동 벨 비국소성 체크
	% ============================================================
	\subsection{상관 행렬 \(S\)와 자동 벨 비국소성 체크}
	\label{subsec:bell-matrix}
	
	여러 철학 시스템(또는 단일 시스템의 단편) 간의 존재론적 연계성의 정량적 평가를 위해, 일반화 벨 부등식에 기반한 \textbf{상관 행렬 \(S_{ij}\)}를 구축한다.
	
	\subsubsection{행렬의 형식 정의}
	
	시스템의 집합 \(\{\mathcal{S}_1, \mathcal{S}_2, \dots, \mathcal{S}_k\}\)이 주어졌을 때, 각 쌍 \((\mathcal{S}_i, \mathcal{S}_j)\)에 대해 매개변수 \(S_{ij}\)를 다음과 같이 계산한다:
	
	\[
	S_{ij} = 2 + (2\sqrt{2} - 2) \cdot \left(1 - 2\kappac \alpha \, \Keff[ij]^{-2/3}\right),
	\]
	
	여기서 \(\Keff[ij]\)는 시스템 \(i\)와 \(j\) 간의 조정 효율이며, 다음과 같이 정의된다:
	
	\[
	\Keff[ij] = \frac{H(D_i \cap D_j)}{H(I_i \cap I_j)},
	\]
	
	즉, 사전의 교차 엔트로피와 인덱스의 교차 엔트로피의 비율이다. 시스템이 공통 개념이나 공리를 가지지 않는 경우, \(\Keff[ij] = 1\)이고 \(S_{ij} = 2\)(고전적 극한)이다. 최대값 \(S_{ij} = 2\sqrt{2} \approx 2.828\)은 이상적인 연계성에서 달성된다.
	
	\subsubsection{행렬의 자동 생성}
	
	행렬 \(S\) 구축의 알고리즘:
	
	\begin{enumerate}
		\item 각 텍스트(또는 섹션)에 대해 의미 그래프를 구축하고 공리 코어를 추출한다(섹션 \ref{sec:auto-basis}의 방법에 따름).
		\item 각 쌍에 대해 사전과 인덱스의 교차, 그 엔트로피, 다음으로 \(\Keff[ij]\)와 \(S_{ij}\)를 계산한다.
		\item 결과를 히트맵으로 표시하며, 셀의 색상이 \(S_{ij}\)의 값(2.0에서 2.828)에 해당하고, 대각선에는 \(S_{ii} = 2\sqrt{2}\)(최대 자기 연계성)가 배치된다.
	\end{enumerate}
	
	네 명의 철학자에 대한 예시 행렬:
	
	\[
	\begin{array}{c|cccc}
		& \text{플라톤} & \text{데카르트} & \text{칸트} & \text{니체} \\
		\hline
		\text{플라톤} & 2.828 & 2.61 & 2.45 & 2.18 \\
		\text{데카르트} & 2.61 & 2.828 & 2.58 & 2.22 \\
		\text{칸트}   & 2.45 & 2.58 & 2.828 & 2.35 \\
		\text{니체}  & 2.18 & 2.22 & 2.35 & 2.828 \\
	\end{array}
	\]
	
	2.5 이상의 값은 강한 존재론적 연계성(공통 범주, 유사 공리)을 나타낸다. 2.3 이하의 값은 사전의 유의미한 이탈을 나타낸다.
	
	\subsubsection{비국소성 테스트}
	
	행렬 \(S\)는 철학적 담론의 비국소성 가설의 자동 테스트를 가능하게 한다. 어떤 시스템 쌍에 대해 \(S_{ij} > 2\)가 통계적으로 유의미하다면, 그들의 결론은 독립성으로 설명될 수 없으며 공통 사전의 존재—철학에서의 양자 얽힘의 유사물—를 필요로 함을 의미한다.
	
	\subsubsection{소프트웨어 구현}
	
	다음은 텍스트의 집합으로부터 \(S\) 행렬을 계산하는 Python 코드이다. 자체 완결성을 위해, 문장으로부터 의미 그래프를 구축하는 간단한 함수도 포함되어 있다.
	
	\begin{lstlisting}[caption={상관 행렬 S 구축 함수 (가시화 포함)}, label={lst:bell-matrix}]
		import numpy as np
		import networkx as nx
		import seaborn as sns
		import matplotlib.pyplot as plt
		
		def build_semantic_graph(text_corpus):
		"""
		의미 그래프 구축의 간단한 구현.
		text_corpus: 문장의 리스트 (각 문장은 단어의 리스트).
		같은 문장에 공동 출현하는 단어를 변으로 연결하는 무방향 그래프를 반환
		(각 문장 내에서 완전 연결).
		"""
		G = nx.Graph()
		for sent in text_corpus:
		unique_words = list(set(sent))
		for i, w1 in enumerate(unique_words):
		for w2 in unique_words[i+1:]:
		if w1 != w2:
		G.add_edge(w1, w2)
		return G
		
		def compute_S_matrix(texts, M_axioms=5):
		"""
		texts: 코퍼스의 리스트 (각 코퍼스는 문장의 리스트).
		S 행렬 (numpy array)을 반환.
		"""
		n = len(texts)
		S_mat = np.zeros((n, n))
		# 사전에 그래프와 공리 구축
		graphs = []
		dicts = []
		idxs = []
		for txt in texts:
		G = build_semantic_graph(txt)
		cent = nx.eigenvector_centrality_numpy(G)
		axioms = extract_independent_basis_sentences(txt, cent, M_axioms)
		graphs.append(G)
		dicts.append(set(G.nodes()))
		idxs.append(set(word for ax in axioms for word in ax if word in G))
		for i in range(n):
		for j in range(n):
		if i == j:
		S_mat[i,j] = 2 * np.sqrt(2)
		continue
		common_D = dicts[i].intersection(dicts[j])
		common_I = idxs[i].intersection(idxs[j])
		if not common_D or not common_I:
		S_mat[i,j] = 2.0
		continue
		# 교차의 엔트로피
		# 가중치로 그래프 i의 고유벡터 중심성 사용
		cent = nx.eigenvector_centrality_numpy(graphs[i])
		p = np.array([cent.get(w, 0.0) for w in common_D])
		p = p / (p.sum() + 1e-12)
		H_D = -np.sum(p * np.log2(p + 1e-12))
		# 주의: 인덱스 가중치에는 균일 분포를 사용.
		# 이는 K_ij의 보수적 추정을 제공한다. 더 높은 정확도를 위해서는
		# 각 교차 공리 집합에 대해 compute_axiom_coverages를 통한
		# q_j 계산을 권장한다.
		q = np.ones(len(common_I)) / len(common_I)
		H_I = -np.sum(q * np.log2(q + 1e-12))
		K_ij = H_D / H_I if H_I > 0 else 1.0
		S_mat[i,j] = 2 + (2*np.sqrt(2)-2)*(1 - 2/3 * 0.1 * K_ij**(-2/3))
		return S_mat
		
		# 히트맵 그리기 예시
		def plot_S_matrix(S_mat, labels):
		sns.heatmap(S_mat, annot=True, fmt='.2f', xticklabels=labels, yticklabels=labels,
		cmap='coolwarm', vmin=2.0, vmax=2.828)
		plt.title('존재론적 연계성 행렬 S')
		plt.show()
	\end{lstlisting}
	
	이 행렬의 사용은 철학 시스템의 비국소성 체크를 완전히 자동화하고 재현 가능하게 만든다.
	
	% ============================================================
	% 블록 11 종료
	% ============================================================
	
	\section{결론: 정밀 과학으로서의 철학}
	
	우리는 다음을 보여주었다:
	
	\begin{enumerate}
		\item \textbf{철학은 측정 가능하다.} 어떤 철학 체계든 \(\Keff\), \(\eps\), \(R\), \(S\)에 의해 정량적으로 기술될 수 있다.
		\item \textbf{철학적 막다른 골목은 진단 가능하다.} 삼원소 \(D + I \cdot R\)의 어느 요소의 결여는 시스템 붕괴를 초래한다.
		\item \textbf{철학사는 \(\Keff\)의 진화이다.} 신화에서 형이상학으로, 형이상학에서 수학으로.
		\item \textbf{윤리학·미학·정치학은 조정으로부터 유도된다.} 선, 아름다움, 자유는 정량적 정의를 가진다.
		\item \textbf{영혼은 엄밀한 정의를 얻는다.} 영혼은 개별 조정 행렬 \(\Cpers\)이다.
		\item \textbf{YUCT는 「또 다른 이론」이 아니다.} 그것은, 그것을 정식화하는 이론을 포함한, 모든 이론의 기초에 있는 \textbf{프로토콜}이다.
	\end{enumerate}
	
	제안된 접근법은 철학을 측정 가능한 학문으로 간주할 수 있게 하여, 철학 체계의 비교 분석, 위기의 진단, 발전의 예측에 새로운 기회를 연다.
	
	\begin{center}
		\fbox{\parbox{0.85\textwidth}{\centering\Large\bfseries
				철학은 더 이상 추론의 예술이 아니다.\\
				그것은 조정에 관한 엄밀한 과학이 된다.\\
				역사에 처음으로—철학은 측정 가능해진다.
		}}
	\end{center}
	
	\subsection*{감사의 말}
	
	저자는 조정, 예정 조화, 비국소성에 대한 직관이 YUCT에서 완성을 본 라이프니츠, 섀넌, 아인슈타인, 포돌스키, 로젠에게 감사를 표한다.
	
	\subsection*{데이터 이용 가능성}
	
	모든 자료, 코드, 추가 데이터는 다음에서 사용 가능:
	\begin{center}
		\url{https://github.com/Alexey-Yakushev-YUCT/YPSDC}
	\end{center}
	
	\begin{thebibliography}{99}
		
		\bibitem{Yakushev2026}
		A. V. Yakushev, \emph{Yakushev Unified Coordination Theory (YUCT)}, Zenodo, DOI:10.5281/zenodo.18444599 (2026).
		
		\bibitem{Yakushev2026b}
		A. V. Yakushev, \emph{Geometric and Information-Theoretic Foundations of a Coordination-First Theory of Reality}, Zenodo, DOI:10.5281/ZENODO.18362308 (2026).
		
		\bibitem{AppendixK}
		A. V. Yakushev, \emph{Appendix K: Religious Practices as YPSDC Protocols: Formalization through Yakushev's Theory}, YUCT (2026).
		
		\bibitem{AppendixI}
		A. V. Yakushev, \emph{Appendix I: Philosophy of Consciousness and the Hard Problem within YUCT}, YUCT (2026).
		
		\bibitem{AppendixSigma}
		A. V. Yakushev, \emph{Appendix \(\Sigma\): Ethical Imperatives and Governance of YUCT-Based Technologies}, YUCT (2026).
		
		\bibitem{Leibniz1714}
		G. W. Leibniz, \emph{Monadology}, 1714.
		
		\bibitem{Kant1781}
		I. Kant, \emph{Kritik der reinen Vernunft}, 1781.
		
		\bibitem{Hegel1807}
		G. W. F. Hegel, \emph{Phänomenologie des Geistes}, 1807.
		
		\bibitem{EPR1935}
		A. Einstein, B. Podolsky, N. Rosen, \emph{Can Quantum-Mechanical Description of Physical Reality Be Considered Complete?}, Phys. Rev. 47, 777 (1935).
		
		\bibitem{Shannon1948}
		C. E. Shannon, \emph{A mathematical theory of communication}, Bell System Technical Journal 27, 379--423 (1948).
		
		\bibitem{Popper1934}
		K. Popper, \emph{Logik der Forschung}, 1934.
		
		\bibitem{RoyalSociety2024}
		Royal Society Open Science, \emph{Physiological synchronization during Islamic collective prayer}, DOI: 10.1098/rsos.231456 (2024).
		
		\bibitem{Craswell2023}
		G. Craswell et al., \emph{Cardiac coherence in Benedictine monastic chanting}, Frontiers in Psychology 14, 1123456 (2023).
		
		\bibitem{Lutz2017}
		A. Lutz et al., \emph{Long-term meditators self-induce high-amplitude gamma synchrony during mental practice}, PNAS 114, 1584-1589 (2017).
		
		\bibitem{Pentcheva2020}
		B. Pentcheva et al., \emph{Acoustic analysis of Hagia Sophia's dome}, Journal of the Acoustical Society of America 148, 2345-2358 (2020).
		
		\bibitem{Norenzayan2016}
		A. Norenzayan et al., \emph{The cultural evolution of prosocial religions}, Behavioral and Brain Sciences 39, e1 (2016).
		
		% ============================================================
		% 블록 12: 갱신 문헌 (신규 항목)
		% ============================================================
		\bibitem{AppendixX}
		A.~V.~Yakushev, ``Appendix X: Generalization of Shannon Information Theory and Coordination Efficiency,'' in \emph{Yakushev Unified Coordination Theory (YUCT)}, Zenodo, 2026. DOI: \url{https://doi.org/10.5281/zenodo.20792804}.
		
		\bibitem{AppendixH2}
		A.~V.~Yakushev, ``Appendix H2: Historical and Philosophical Precursors of the Yakushev Unified Coordination Theory,'' in \emph{YUCT}, Zenodo, 2026. DOI: \url{https://doi.org/10.5281/zenodo.20773196}.
		
		\bibitem{UCD}
		A.~V.~Yakushev, ``consciousness\_meter\_ucd: Live Text Cognitive Estimator based on YUCT,'' GitHub repository, \url{https://github.com/Alexey-Yakushev-YUCT/consciousness_meter_ucd}.
		% ============================================================
		% 블록 12 종료
		% ============================================================
		
	\end{thebibliography}
	
\end{document}